% wheat-sourdough-ar.tex — ch. 9 الخبز الحامض القمحي, The Sourdough
% Framework: Arabic Edition (R2 unit 18)
%
% Original: The Sourdough Framework by Hendrik Kleinwächter,
% the-sourdough-framework.com, CC BY-SA 4.0. The Arabic edition inherits
% CC BY-SA 4.0.
%
% RECONSTRUCTION-FIRST (Phase 28R): transcribed from the accepted edition
% text layer (book-ar-fixed.pdf, b78c2a28…, PRIMARY — accepted pp. 72–117),
% cross-checked against the recovery reference
% (recovery-r0/text/09-wheat-sourdough.txt, phase-23 make4ht conversion of
% the accepted Arabic sources); upstream book/wheat-sourdough/
% wheat-sourdough.tex used for STRUCTURE ONLY (float environments, labels,
% cites, footnote order, subsection short titles). No retranslation.
% Structure: \chapter + quoting intro; 14 sections (9.1–9.14), 6
% subsections in 9.12, 2 in 9.13; 7 flowcharts (fig-starter-readiness-ar
% shared with ch. 5), 4 tables (tables-ar/), ~24 photos carried over
% unchanged; 14 footnotes; 4 cites. Float placement specs follow the
% accepted-edition-warranted documented deviation [!htb] → [htb] (unit 12
% convention). Labels verbatim from upstream. Latin runs wrapped
% \textenglish{} (unit 11 sans-context routing); ranges
% \LR{\textenglish{N\textendash M}} (unit 9 dash-font convention);
% \qty rides \LR{} islands. Section titles follow the accepted edition
% (9.2 تهيئة البادئة, 9.5 كم مقدار البادئة؟, 9.6 أوتولايز, 9.7
% الفيرمنتوليز — the recovery TOC's forms, not upstream's literal
% translations).

\chapter{الخبز الحامض القمحي}%
\label{ch:wheat-sourdough}

\begin{quoting}
في هذا الفصل ستتعلّم كيفية صنع الخبز الحامض القمحي القائم بذاته.
\end{quoting}

\begin{figure}[htb]
  \centering
  \includegraphics[width=\textwidth]{loaf-pan-free-standing}
  \caption[رغيف قائم بذاته ورغيف في قالب]{رغيف خبز حامض قائم بذاته بجانب رغيف مصنوع في قالب الخبز. يَعُدُّ كثير من الخبازين الخبزَ الحامض القائم بذاته التحدّيَ الأعلى في فنّ خبز الحامض.}
\end{figure}

الخبز الحامض القائم بذاته هو نوع الخبز المفضّل عندي. فهو يجمع قشرة مقرمشة رائعة ونكهة ممتازة ولُبًّا طريًّا هشًّا. إنه الخبز الذي يلتهمه أصدقائي وعائلتي على الفور. ولسوء الحظ، فإن صنع هذا النوع من الخبز يتطلّب جهدًا وصبرًا ومهارةً أكبر بكثير من الأنواع الأخرى. فأنت مضطرّ إلى موازنة عملية التخمير بدقّة تامّة: لا يمكنك أن تخمّر لمدة قصيرة أكثر من اللازم، ولا لمدة طويلة أكثر من اللازم أيضًا. وتتطلّب التقنيات التي ستحتاج إلى تعلّمها قدرًا إضافيًّا من المهارة. احتجتُ بنفسي إلى عدة محاولات حتى أتقنتُها، وواجهتُ تحدّيات عدة: كان عندي دقيق غير مناسب، ولم أكن أعرف كيف أستخدم فرني كما يجب، ومتى ينبغي أن أوقف التخمير؟ المعلومات المنتشرة حول هذا الموضوع كثيرة جدًّا؛ فقدّعتُ في معظمها وجرّبتُ كل شيء تقريبًا. وكثيرًا ما كانت المعلومات خاطئة، وفي أحيان أخرى كنت أكتشف قطعة من الألغاز قيّمة جديدة. كان جمع هذه المعلومات وفكّ شفرة بعضها إحدى دوافعي الرئيسية لبدء \textenglish{\texttt{The Bread Code}}. وخلاصة ما تعلمته أنه لا وصفة يمكن أن تتبعها اتباعًا أعمى؛ فستضطرّ دومًا إلى تكييف الوصفة مع الأدوات والبيئة المتاحة لديك محليًّا.

لكن لا تقلق. فبعد قراءة هذا الفصل ستعرف كل العلامات التي ينبغي أن تراقبها، وستصبح قادرًا على قراءة عجينك، وستتحوّل إلى خباز هاوٍ واثق من نفسه قادر على خَبز الخبز في بيته، وفي الأماكن المرتفعة والمنخفضة على السواء، صيفًا وشتاءً، في بيت صديق له، بل حتى أثناء الإجازة. علاوة على ذلك ستعرف كيف توسّع إنتاجك من رغيف واحد إلى مئة رغيف. فإذا رغبت يومًا في فتح مخبز، فاعتبر هذه المعرفة أساسك الذي تبنِي عليه.

وإتقانك لهذه العملية سيمكّنك من صنع خبز مذهل النكهة يفوق أي خبز يُشترى من المتجر.

\section{العملية}

\begin{flowchart}[htb]
% fig-wheat-sourdough-process-ar.tex — Arabic edition of upstream
% figures/fig-wheat-sourdough-process.tex (ch. 9, المخطط الانسيابي 9.1)
%
% Original: The Sourdough Framework by Hendrik Kleinwächter,
% the-sourdough-framework.com, CC BY-SA 4.0. The Arabic edition inherits
% CC BY-SA 4.0.
%
% R1-FIGURES route: re-derived from the upstream source
% figures/fig-wheat-sourdough-process.tex (node styles from
% figures/flowcharts_tikz.tex and the node/edge geometry verbatim), with
% the label text recovered from the accepted edition's render of p. 73
% (book-ar-fixed.pdf, b78c2a28…) — label recovery against a visual
% reference, not retranslation. Line splits inside nodes follow the
% accepted edition's render exactly (اخلط/المكونات، ابِن قوة/العجين،
% هل تخبز/رغيفًا واحدًا؟). Arabic label text is wrapped in
% \textarabic{}; digit/Latin runs ride \LR{} islands (Ruling 1
% uniform-Pagella figure face, ACCEPTED 2026-11-15). The decision edge
% labels نعم/لا follow the accepted render.

\begin{tikzpicture}[node distance = 3.2cm, auto]
  \node [start] (init) {\textarabic{بادئة جاهزة}};
  \node [block, right of=init] (mix_ingredients) {\textarabic{اخلط}\\ \textarabic{المكونات}};
  \node [block, right of=mix_ingredients] (dough_strength) {\textarabic{ابِن قوة}\\ \textarabic{العجين}};
  \node [block, right of=dough_strength] (bulk) {\textarabic{التخمير الأول}};
  \node [decision, below of=bulk] (divide_test) {\textarabic{هل تخبز}\\ \textarabic{رغيفًا واحدًا؟}};
  \node [block, right of=divide_test] (divide) {\textarabic{قسّم}};
  \node [block, below of=divide] (preshape) {\textarabic{التشكيل المسبق}};
  \node [block, below of=divide_test] (shape) {\textarabic{شكّل}};
  \node [block, left of=shape] (proof) {\textarabic{التخمير النهائي}};
  \node [success, left of=proof] (bake) {\textarabic{اخبز}};
  \path [line] (init) -- (mix_ingredients);
  \path [line] (mix_ingredients) -- (dough_strength);
  \path [line] (dough_strength) -- (bulk);
  \path [line] (bulk) -- (divide_test);
  \path [line] (divide_test) -- node{\textarabic{نعم}} (shape);
  \path [line] (divide_test) -- node{\textarabic{لا}} (divide);
  \path [line] (divide) -- (preshape);
  \path [line] (preshape) -- (shape);
  \path [line] (shape) -- (proof);
  \path [line] (proof) -- (bake);
\end{tikzpicture}

\caption[العملية النموذجية لصنع خبز حامض قمحي]{العملية النموذجية لصنع خبز حامض قمحي.}
\label{fig:wheat-sourdough-process}
\end{flowchart}

تبدأ العملية الكاملة لصنع خبز حامض رائع بتهيئة بادئتك الحامضة. ومفتاح إتقان هذه العملية هو إدارة التخمير كما يجب، وأساس ذلك أن يكون لديك بادئة حامضة نشطة ومعافاة.

وبعد أن تجهز بادئتك، تنتقل إلى خلط جميع المكونات. تريد أن تجنّس بادئتك الحامضة جيدًا، وبذلك تضمن تخميرًا متساويًا في العجين بأكمله.

وبعد استراحة قصيرة، تنتقل إلى بناء قوة العجين. فالعجن يخلق شبكة غلوتين قوية، وهذا أمر جوهري لحصر ثاني أكسيد الكربون \LR{CO\textsubscript{2}} الناتج عن التخمير حصرًا صحيحًا.

وبعد العجن يبدأ التخمير الأول. ويُسمّى تخميرًا أولًا لأنك تخمّر عادةً عدة أرغفة معًا في كتلة واحدة. وسيتطلب فهم متى توقف هذه الخطوة بعض التدريب، لكن لا شيء يدعو للقلق: ستتعلم العلامات الدقيقة التي تبحث عنها.

وبعد اكتمال هذه المرحلة تحتاج إلى تقسيم كتلة العجين الكبيرة إلى قطع أصغر، وتشكيل كل قطعة تشكيلًا مسبقًا. فهذا يتيح لك تطبيق مزيد من قوة العجين وتشكيل أرغفة أكثر انتظامًا.

ثم تأتي مرحلة التخمير النهائي حيث تُكمل عملية التخمير. وبحسب وقتك يمكنك تخميره في درجة حرارة الغرفة أو في الثلاجة. وإتقان التخمير النهائي هو ما يحوّل رغيفك الجيد إلى رغيف عظيم.

وأخيرًا تختم العملية كلها بالخَبز، حيث ستتعلم خيارات مختلفة لتوليد البخار لعجينك كما يجب، فيحصل عجينك على انتفاخ جميل في الفرن. وفي المرحلة الثانية من عملية الخَبز تكمل بناء القشرة.

وتعتمد كل هذه الخطوات على بعضها بعضًا؛ فلن تحصل على الخبز المثالي إلا إذا أتقنت كل خطوة منها.

\section{تهيئة البادئة}%
\label{sec:readying-starter}

أهم جزء في عملية صنع الخبز هو بادئتك. فالبادئة هي ما يبدأ التخمير في عجينك الرئيسي، وإذا كانت بادئتك غير سليمة فسيعاني عجينك الرئيسي المشكلات نفسها أثناء التخمير. فخصائص بادئتك تنتقل إلى عجينك الرئيسي: إذا لم تتوازن في بادئتك نسبة الخميرة إلى البكتيريا فلن تتوازن في عجينك الرئيسي أيضًا.

\begin{flowchart}[htb]
% fig-starter-readiness-ar.tex — Arabic edition of the sourdough starter
% readiness flowchart (ch. 5, المخطط الانسيابي 5.2)
%
% Original: The Sourdough Framework by Hendrik Kleinwächter,
% the-sourdough-framework.com, CC BY-SA 4.0. The Arabic edition inherits
% CC BY-SA 4.0.
%
% R1-FIGURES route: re-derived from the upstream post-#492 source
% figures/fig-starter-readiness.tex (node styles from
% figures/flowcharts_tikz.tex and the node/edge geometry verbatim,
% including the decision_start→feed_yes_branch.north west and the
% size_check/smell_check south anchors), with the label text recovered
% from the accepted edition's render of p. 40 (book-ar-fixed.pdf,
% b78c2a28…) — label recovery against a visual reference, not
% retranslation. Line splits inside nodes follow the accepted edition's
% render exactly. The upstream r@{}l tabulars in the ratio nodes are set
% as plain \\ lines: the accepted edition's render shows four plain
% lines per ratio node. Arabic label text is wrapped in \textarabic{};
% digit/Latin runs ride \LR{} islands (Ruling 1 uniform-Pagella figure
% face, ACCEPTED 2026-11-15). The two feed-duration nodes render the
% range with a single trailing unit («6–12 h»), as the accepted render's
% per-char positions show; the caption keeps the repeated-unit
% \qtyrange form («6 h–12 h») that the accepted caption shows.

\begin{tikzpicture}[node distance = 3cm, auto]
  \node [start] (init) {\textarabic{أنشئ}\\ \textarabic{بادئة}};
  \node [decision_start, right of=init, node distance=3.5cm] (decision_start) {\textarabic{هل غُّذِيت}\\ \textarabic{البادئة}\\ \textarabic{خلال \LR{3} أيام؟}};
  \node [block, right of=decision_start, text width=7em, node distance=4cm] (feed_no_branch)
      {\textarabic{غّذِ البادئة مرتين:}\\ \textarabic{قبل \LR{\qty{48}{\hour}}}\\ \textarabic{وقبل \LR{\numrange{6}{12}~\unit{\hour}}}};
  \node [block, below of=feed_no_branch, text width=7em, node distance=2.7cm] (feed_yes_branch)
      {\textarabic{غّذِ البادئة}\\ \textarabic{بمدة \LR{\numrange{6}{12}~\unit{\hour}}}\\ \textarabic{قبل صنع العجين.}};
  \node [block, right of=feed_no_branch, text width=7em, node distance=4cm] (high_ratio)
    {\textarabic{اعتمد نسبة \LR{1:10:10}:}\\ \textarabic{\LR{10 g} بادئة،}\\ \textarabic{\LR{100 g} دقيق،}\\ \textarabic{\LR{100 g} ماء.}};
  \node [block, right of=feed_yes_branch, text width=7em, node distance=4cm] (low_ratio)
  {\textarabic{اعتمد نسبة \LR{1:5:5}:}\\ \textarabic{\LR{10 g} بادئة،}\\ \textarabic{\LR{50 g} دقيق،}\\ \textarabic{\LR{50 g} ماء.}};
  \node [decision, below of=high_ratio, node distance=6cm] (size_check)
    {\textarabic{فقاعات؟}\\ \textarabic{وزيادة في الحجم؟}};
  \node [decision, below of=decision_start, node distance=6cm] (smell_check)
    {\textarabic{رائحة خليكية}\\ \textarabic{أم لبنية زبادية؟}};
  \node [success, below of=init, node distance=6cm] (make_dough)
    {\textarabic{جهِّز العجين}};

  \path [line] (init) -- (decision_start);
  \path [line] (decision_start) -- node {\textarabic{لا}} (feed_no_branch);
  \path [line] (decision_start) -- node[below=2pt]{\textarabic{نعم}} (feed_yes_branch.north west);
  \path [line] (feed_yes_branch) -- (low_ratio);
  \path [line] (feed_no_branch) -- (high_ratio);
  \path [line] (high_ratio) -- node[anchor=east, above=2pt] {} ++(2.2,0) |-(size_check);
  \path [line] (low_ratio) -- (size_check);
  \path [line] (size_check) -- node {\textarabic{لا}} (feed_yes_branch.south east);
  \path [line] (size_check) -- node {\textarabic{نعم}} (smell_check);
  \path [line] (smell_check) -- node {\textarabic{لا}} (feed_yes_branch.south west);
  \path [line] (smell_check) -- node {\textarabic{نعم}} (make_dough);

  % braces
  \draw[BC]   (size_check.south) --
      node[below=1em]{\textarabic{تحّقق من جهوزية البادئة قبل الاستخدام}}(smell_check.south);
\end{tikzpicture}

\caption[عملية تهيئة البادئة قبل الخَبز]{العملية التي تفحص بها بادئتك الحامضة عند صنع العجائن القمحية. عمليًّا، أستخدم كثيرًا بادئة حامضة صلبة، والبادئة الصلبة تتميز بنشاط خميرة معزّز، وفي هذه الحالة يمكنك استخدام النسب نفسها المبيّنة في المخطط ما عدا كمية الماء: فالترطيب في البادئة الصلبة يبلغ \LR{\textenglish{50\,\%\textendash 60\,\%}}. إذن ستستخدم نصف كميات الماء المبيّنة، أي إذا أظهر المخطط \LR{100 g} من الماء فستستخدم \LR{\textenglish{50 g\textendash 60 g}} من الماء لبادئتك الصلبة.}
\label{fig:process-starter-wheat-sourdough}
\end{flowchart}

فكّر عمومًا في العجين الذي تخلطه بوصفه بادئة كبيرة مضافة إليها الملح. فبعد خلط جميع المكونات تعود لديك بيئة مفتوحة للمنافسة من جديد، فتبدأ الخميرة والبكتيريا التقاتل مجددًا ليطغى كل منهما على الآخر. فالغذاء متاح بوفرة، وكلها تبذل قصارى جهدها للفوز. وبحسب البادئة التي تخلطها في عجينك قد يطغى بعض الكائنات الدقيقة على غيرها.

الخيار الأول لتحقيق توازن جيد هو تطبيق تغذيات متكررة. فإذا لم تُغذَّ بادئتك منذ مدة طويلة فستطغى البكتيريا، ويحدث هذا إذا تركتَ بادئتك مثلًا غير مستخدمة في الثلاجة. فكلما تراكمت الحموضة أصبحت البيئة معادية للخميرة أكثر وأكثر، بينما تتحمّل بكتيريا حمض اللبنيك هذه البيئة أفضل. وسيكون تخمير عجينك بهذه البادئة أقرب إلى الجانب البكتيري. ومن تطبيق بعض التغذيات تنشط الخميرة من جديد، وكلما قدمت بادئتك زادت مقاومة الخميرة للحموضة. في البداية اضطررتُ إلى تغذية بادئتي من مرتين إلى ثلاث مرات \LR{\textenglish{2\textendash 3}} لإصلاح التوازن، أما مع بادئتي الأكثر نضجًا فتغذية واحدة تكفي لموازنة الكائنات الدقيقة.

يستخدم بعض الناس نسبة \LR{1:1:1} لإنعاش البادئة، وهي جزء واحد من البادئة القديمة (\LR{10 g} مثلًا) وجزء من الدقيق وجزء من الماء. وأرى أن هذا هراء تامّ. فكما ذكرنا، بادئتك عجين مصغّر، ولن تفكر يومًا في نسبة \LR{1:1:1} لصنع عجين. فأنت لن تستخدم أكثر من \LR{20\,\%} من البادئة لصنع العجين. لهذا أنصح باستخدام نسبة \LR{1:5:5} أو نسبة \LR{1:10:10} بحسب نضج بادئتك. وبما أنني أستخدم دائمًا تقريبًا بادئة حامضة أكثر صلابة لأسباب تخمير الخميرة المعزّز لديها (انظر القسم~\ref{sec:stiff-starter}) فإن نسبتي ليست \LR{1:5:5} أبدًا، بل \LR{1:5:2.5} (جزء واحد من البادئة القديمة، وخمسة أجزاء من الدقيق، و\LR{2.5} جزء من الماء). وإذا كان الجو حارًّا جدًّا في مكان سكنك فيمكنك اختيار النسبة \LR{1:10:5} أو \LR{1:20:10} المذكورة سابقًا، وبذلك تُبطئ نضج بادئتك. ويمكنك أيضًا استخدام هذه الحيلة لتوفيق تغذية البادئة مع جدولك: فإذا كانت بادئتك تجهز عادةً في \LR{6} ساعات لكنك تحتاجها اليوم لاحقًا، فما عليك إلا زيادة كمية الدقيق والماء التي تغذّي بها بادئتك. هذه كلها قيم تحتاج إلى التجربة فيها بنفسك، فكل بادئة فريدة وقد تتصرف قليلًا على نحو مختلف.

الخيار الثاني المتاح لك هو كمية البادئة التي تستخدمها لصنع العجين. فكما ذكرنا سابقًا، تعود بادئتك إلى النمو داخل عجينك الرئيسي. وبينما أستخدم عادةً ما يتراوح بين \LR{\textenglish{10\,\%\textendash 20\,\%}} من البادئة بالنسبة إلى الدقيق، أحيانًا أنزل إلى \LR{1\,\%} من البادئة فقط، وبذلك يتوفر للكائنات الدقيقة مجال أوسع لموازنة نفسها أثناء تخمير العجين. فإذا لم أُغذِّ بادئتي الحامضة منذ يوم، فقد أستخدم \LR{5\,\%} من البادئة لصنع العجين. وإذا مدّتُ ذلك إلى يومين بلا تغذيات خفضتُ كمية البادئة أكثر، واخترتُ \LR{1\,\%} من البادئة المذكورة سابقًا. فإذا كان الغذاء شحيحًا جدًّا ستنبّس الكائنات الدقيقة (التبوغ)، وستحتاج إلى النمو من جديد من الجراثيم التي أنتجتها، وفي حالة السبات هذه يستغرق استعادة نشاطها الكامل وقتًا أطول. جرّبتُ عدة مرات صنع العجين مباشرةً من بادئة جافة، ولم أنجح لأن التخمير استغرق وقتًا طويلًا جدًّا: فقد اضطرت الكائنات الدقيقة إلى النمو من الجراثيم ثم بدء التخمير. وكما شرحنا سابقًا هناك حدّ لأزمنة التخمير لأن عجينك يتفكك طبيعيًّا مع الوقت. علاوة على ذلك، تريد أن تتفوق كائناتك الدقيقة على مسبّبات الأمراض الأخرى الموجودة في الدقيق؛ فكلما قلّت كمية البادئة التي تستخدمها أسهل عليها التكاثر. والبادئة القوية ستطغى على الجراثيم الأخرى. وبما أن طريقة تخفيض البادئة تعمل، فإنني أنصح بالخيار الأول أكثر، فهو ينتج خبزًا أفضل بشكل موثوق. أما الخيار الثاني فهو ما ألجأ إليه عادةً حين أغذّي بادئتي في الصباح لكني أفشل في صنع عجين في المساء: لا أريد تغذية بادئتي من جديد في الصباح التالي، ويفضّل أن أصنع العجين مباشرةً دون انتظار، فأستخدم بادئة ناضجة جدًّا وبكمية أقل.

مع الوقت ستألف بادئتك وطريقة تصرفها، وستصبح قادرًا على قراءة علامات نشاطها والحكم على حالتها.

\section{المكونات}

كل ما تحتاجه لصنع خبز حامض رائع هو الدقيق والماء والملح. ويمكنك طبعًا إضافة أشياء أخرى إلى عجينك مثل البذور. وأنا شخصيًّا أحب النكهة القوية للقمح كامل الحبة، لذا أحب إضافة ما يتراوح بين \LR{\textenglish{20\,\%\textendash 30\,\%}} من دقيق القمح كامل الحبة إلى الخليط. ويمكنك أيضًا صنع هذه الوصفة بدقيق القمح كامل الحبة مباشرةً بنسبة \LR{100\,\%}، وفي هذه الحالة ابحث عن دقيق قمح كامل الحبة قوي مصنوع من دقيق ذي بروتين أعلى. وإذا لم تحب القمح كامل الحبة فيمكنك حذفه من الوصفة، واستبدال الكمية المذكورة ببساطة بدقيق الخبز. أما ما ينبغي الانتباه إليه في دقيق القمح كامل الحبة فهو نشاطه الإنزيمي الأعلى: فبإضافة بعض دقيق القمح كامل الحبة ستسرّع عملية التخمير كلها.

وخاصةً في بداياتك أنصح باستخدام دقيق الخبز الذي يحتوي على غلوتين أكثر من الدقيق متعدد الاستخدامات أو دقيق الكعك، وهذا جوهري حين تحاول خَبز رغيف قائم بذاته بالبادئة الحامضة.

إليك في الأسفل وصفة نموذجية لرغيف واحد مع حسابها بالحساب الخبازي:

\begin{itemize}
  \item \LR{400 g} من دقيق الخبز
  \item \LR{100 g} من دقيق القمح كامل الحبة
  % Manual unit so we can use emphasis
  \item \emph{المجموع: \LR{500 غ} من الدقيق}
  \item \LR{\textenglish{300 g\textendash 450 g}} من ماء بدرجة حرارة الغرفة (من \LR{60\,\%} إلى \LR{90\,\%}). المزيد عن هذا الموضوع في الفصل التالي.
  \item \LR{50 g} من البادئة الحامضة الصلبة (\LR{10\,\%})
  \item \LR{10 g} من الملح (\LR{2\,\%})
\end{itemize}

وإذا أردت صنع خبز أكثر فما عليك إلا زيادة الكميات بحسب كمية الدقيق المتاحة لديك. لنقل أن لديك \LR{2000 g} من الدقيق، ستبدو الوصفة هكذا:

\begin{itemize}
  \item \LR{1600 g} من دقيق الخبز
  \item \LR{400 g} من دقيق القمح كامل الحبة
  % Manual unit so we can use emphasis again
  \item \emph{المجموع: \LR{2000 غ} من الدقيق}، أي ما يكفي \LR{4} أرغفة
  \item من \LR{1200 g} إلى \LR{1800 g} من ماء بدرجة حرارة الغرفة (من \LR{60} إلى \LR{90\,\%})
  \item \LR{200 g} من البادئة الحامضة الصلبة (\LR{10\,\%})
  \item \LR{40 g} من الملح (\LR{2\,\%})
\end{itemize}

هذا هو جمال الحساب الخبازي: أعد حساب النسب المئوية فحسب، وأنت جاهز. وإذا لم تكن متأكدًا من كيفية عمل ذلك، فاطلع على القسم~\ref{sec:bakers-math} الكامل الذي يتناول الموضوع بالتفصيل.

\section{الترطيب}

يعني الترطيب مقدار الماء الذي تستخدمه بالنسبة إلى دقيقك. فحين بدأتُ صنع الخبز كنت أخطئ في هذا دائمًا: كنت أتبع وصفة من الإنترنت، ولم يبدُ عجيني أبدًا مثل العجين الظاهر في الوصفة. فكمية الماء التي يحتاجها دقيقك ليست ثابتة، بل تعتمد على الدقيق الذي بين يديك.

فحين تتلامس الحبّة مع الماء أول مرة، تمتص الطبقات الخارجية الماء، ولهذا حين تستخدم القمح كامل الحبة (الذي ما يزال يحتوي هذه الطبقات) عليك استخدام قليل من الماء الإضافي.

وبتكوين خيوط الغلوتين يُمتصّ الماء في مصفوفة الغلوتين في عجينك، وكلما ارتفعت قيمة البروتين أمكن استخدام ماء أكثر.

يحب بعض الخبازين استخدام عجائن مرتفعة الترطيب لصنع خبز أكثر هشاشة وفتوحة.\footnote{أحيانًا تكاد تبدو المقارنة بين الخبازين على أنها مقارنة للبراعة: كلما تمَكَّن الخباز من التعامل مع ماء أكثر اعتُبر أكثر مهارة.} وسبب ذلك هو قابلية العجين المحدّسة الأفضل: فكلما كان العجين رطبًا أكثر كان مدّه أسهل؛ فحين تشدّه يحافظ على شكله. وفي المقابل يحافظ العجين الصلب جدًّا (منخفض الترطيب) على شكله مدة أطول. وللتصور، فكّر في عجينك القابل للمد بوصفه بالونًا، والعجين الصلب مثل إطار سيارة: فالخميرة تجد صعوبة أكبر بكثير في نفخ إطار السيارة مقارنةً بالبالون، لأن مطاط إطار السيارة أقل قابلية للمد بكثير، ويحتاج نفخه إلى قوة أكبر بكثير. لهذا السبب سينتفخ العجين القابل للمد أكثر في الفرن، فيبدو الرغيف أكبر بصريًّا ويقدّم بنية لُبّ أكثر فتوحة وانفتاحًا.

ورغم أن هذا قد يبدو رائعًا، فإن الترطيب العالي يسبّب آثارًا جانبية عدة:

\begin{enumerate}
  \item يصبح عجينك أصعب في التعامل، وسيكون أكثر التصاقًا.
  \item يلزم عجن عجينك مدة أطول لبناء شبكة غلوتين سليمة.
  \item أثناء التخمير قد يصبح عجينك قابلًا للمد أكثر من اللازم فيفقد بعض قوة العجين، ولتجاوز ذلك تُطبَّق الطيّات والشدّات مقارنةً بالعجين العادي، ما يتطلب منك استثمار عمل أكبر بكثير.
  \item يصبح التشكيل أكثر إزعاجًا بكثير لأن العجين لزج جدًّا.
  \item يمكن أن يلتصق العجين بسلة التخمير بسهولة أكبر أثناء التخمير النهائي.
  \item إذا انتظرت طويلًا جدًّا أثناء التخمير النهائي فلن يبقى للعجين قوة كافية ليشدّ إلى الأعلى وسيبقى مسطحًا.
  \item عمومًا، كلما ارتفع محتوى الماء زاد التخمير البكتيري؛ فالعجين الأكثر رطوبة يفكّك الغلوتين أسرع من العجين الأصلب، ولهذا عليك أن تبدأ التخمير ببادئة حامضة في حالة مثالية. ويستخدم الخبازون عملية تُسمّى أوتولايز (autolysis) لتقصير زمن التخمير الرئيسي وتجاوز هذه المشكلة.
  \item وقد يُستشعر اللُّبّ في النهاية لزجًا إلى حدٍّ ما، لأنه ما يزال يحتوي كمية كبيرة من الماء. أنا أحب هذا القوام الرطب، لكن الأمر يبقى مسألة ذوق شخصي.
\end{enumerate}

ولتحقيق عجين مرتفع الترطيب من الأفضل أن تضيف الماء إلى عجينك ببطء: ابدأ بترطيب \LR{60\,\%} ثم أضف قليلًا من الماء ببطء، وعجن من جديد حتى يُمتصّ الماء، وكرّر وأضف ماءً أكثر. وبما أن عجينك شكّل شبكة غلوتين بالفعل، يمكن امتصاص الماء الجديد بسهولة أكبر بكثير، وستندهش من كمية الماء التي يستطيع عجينك امتصاصها. تُعرف هذه الطريقة عمومًا بطريقة الباسيناج (bassinage)، وسنعود إليها لاحقًا. وباعتماد هذه التقنية استطعتُ بسهولة أن أرفع دقيقًا منخفض الغلوتين إلى ترطيب \LR{80\,\%}. وهذه أيضًا طريقتي المفضّلة عند صنع العجين الآن: أستمر في إضافة الماء حتى أشعر بأن العجين وصل إلى القوام المناسب. وكلما خبزت خبزًا أكثر طوّرت إحساسًا وعينًا أفضل لعجينك. وقد يكون هذا مرهقًا بعض الشيء عند الخلط باليد، وهو أسهل كثيرًا عند استخدام عجّنة كهربائية.

وعلى أي حال، فإن رفع الترطيب يتطلب محاولات وأخطاء كثيرة. لكن هناك خيارًا واحدًا يجعل الأمور أسهل ويسبّب صداعًا أقل: التخمير البطيء. فتحصل على مزايا قابلية المد نفسها التي يقدّمها الترطيب العالي بمجرد أن تترك عجينك يتخمّر مدة أطول. وإبطاء عملية التخمير أمر سهل: استخدم بادئة حامضة أقل أو خمّر في بيئة أبرد.

وللأسباب التالية تكون للتخمير البطيء مزاياه. فكما شرحنا سابقًا، فإن إنزيم البروتييز والبكتيريا معًا يفكّكان شبكة الغلوتين في عجينك؛ فكلما تقدّم التخمير أصبح عجينك قابلًا للمد تلقائيًّا، لأن طبقات المطاط في إطار سيارتك تتحول وتُؤكل ببطء، حتى يتحوّل إطار سيارتك في النهاية إلى بالون يمكن نفخه بسهولة بالغة. وإذا انتظرت طويلًا جدًّا سينفجر البالون: لن يبقى لديك غلوتين، وسيصبح عجينك لزجًا جدًّا. والعثور على النقطة المثلى بين أكل المطاط بما يكفي وبين الإفراط فيه هو جوهر الخبز الحامض القمحي المثالي. لكن لا تقلق؛ فبعد قراءة هذا الفصل ستكون الأدوات المناسبة بين يديك.

ويمكن ملاحظة مزايا التخمير البطيء جيدًا عند التجربة بعجين خميرة سريع التخمير (\LR{1\,\%} خميرة جافة بالنسبة إلى الدقيق). فلُبّ مثل هذا العجين لا يكون أبدًا منفتحًا مثل عجين مصنوع بالبادئة الحامضة. علاوة على ذلك، لا يستطيع إنزيم البروتييز أداء عمله في مدة تخمير قصيرة هذه. وتضيف المخابز الصناعية الكبيرة مالتًا نشطًا يحتوي إنزيمات أكثر بكثير، وبذلك تقصّر الوقت اللازم لصنع العجين. وستجد المالت غالبًا ضمن مكونات خبز المتاجر، وهذه حيلة ذكية، لكن خبز التخمير السريع الفائق هذا سيتميز بلُبّ كثيف نسبيًّا وليس هشًّا، لأن الغلوتين القليل جدًّا فقط يتفكك عند إتمام زمن التخمير في \LR{1} ساعة واحدة. ولو أبطأت الأمور لبدو العجين مختلفًا تمامًا. جرّب ذلك مرة أخرى بخميرة أقل بكثير؛ فهذا هو سرّ البيتزا النابوليتانية: كمية ضئيلة جدًّا من الخميرة فقط تُستخدم لصنع العجين. وصفتي الافتراضية للبيتزا تتطلب نحو \LR{150 mg} من الخميرة الجافة لكل \LR{kg} من الدقيق. جرّب بنفسك في المرة القادمة التي تصنع فيها عجين خميرة: حاول دفع التخمير إلى \LR{8} ساعات على الأقل. سيكون الفرق مذهلًا: ستصنع خبزًا بلُبّ أكثر هشاشة وفتوحة بكثير، وستتحسن نكهة العجين تحسنًا جذريًّا، وستصبح قشرتك أكثر قرمشة وبنكهة أفضل، لأن الأميلازات قد حوّلت النشويات إلى سكريات أبسط تعمل على التحبّب بشكل أفضل أثناء الخَبز. وإذا لم تتعلم من هذا الكتاب سوى شيء واحد، فليكن أن التخمير البطيء هو مفتاح صنع الخبز الرائع.

لهذا السبب، ترطيبي الافتراضي أقل بكثير من ترطيب خبازين آخرين؛ فأنا أفضّل التخمير البطيء في وصفاتي. والنقطة المثلى لدقيق عندي تقع حول ترطيب \LR{70\,\%}. وأكرر: هذه قيمة ذاتية جدًّا تناسب دقيق أنا.

وإذا كنت تبدأ للتوّ بدفعة جديدة من الدقيق، فأنصحك بإجراء الاختبار التالي، فهو سيساعدك على تحديد النقطة المثلى لقدرات ترطيب دقيقك.

اصنع خمسة أوانٍ، كل واحد منها بـ\LR{100 g} من الدقيق، وأضف كميات ماء متزايدة قليلًا إلى كل واحد منها:

\begin{itemize}
  \item \LR{100 g} من الدقيق و\LR{55 g} من الماء
  \item \LR{100 g} من الدقيق و\LR{60 g} من الماء
  \item \LR{100 g} من الدقيق و\LR{65 g} من الماء
  \item \LR{100 g} من الدقيق و\LR{70 g} من الماء
  \item \LR{100 g} من الدقيق و\LR{75 g} من الماء
\end{itemize}

ثم امزج خليط الدقيق والماء حتى ترى أنه لم يبق منه أي كتل دقيق. انتظر \LR{15} دقيقة وعُد إلى عجينك، وافصل العجين بعناية بين يديك؛ ينبغي أن يكون عجينك مرنًا ومتماسكًا تمامًا. مدّ عجينك حتى يصبح رقيقًا جدًّا ثم ارفعه مقابل ضوء: ينبغي أن تستطيع الرؤية من خلاله. وخلط الدقيق والماء الذي ينكسر دون أن ترى نافذة (window pane) هو منطقة ممنوعة لك. فاختر عجينًا بترطيب أقل من هذه القيمة؛ فستعرف مثلًا أن مزيج دقيقك يصل إلى ترطيب \LR{65\,\%}. واستخدم بقايا هذه التجربة لتغذية بادئتك.

\begin{figure}[htb]
  \centering
  \includegraphics[width=\textwidth]{window-pane-effect}
  \caption[اختبار النافذة]{يتيح لك اختبار النافذة أن ترى بنفسك ما إذا كنت قد طوّرت الغلوتين تطويرًا كافيًا.}
\end{figure}

ومن الناحية الاقتصادية، يعدّ الماء أرخص مكونات عجين الخبز. ففي المخبزة يزن العجين الأكثر ترطيبًا وزنًا أكبر وتبلغ تكاليف إنتاجه أقل، فيرتفع الربح. لكن هذا يأتي مقابل زيادة تكاليف العمل وزيادة احتمالات الفشل بسبب صعوبة التعامل المتزايدة.

\section{كم مقدار البادئة؟}

يستخدم معظم الخبازين نحو \LR{20\,\%} من البادئة الحامضة بالنسبة إلى وزن الدقيق. وأنا أنصح بالانخفاض أكثر بكثير، إلى نحو \LR{\textenglish{5\,\%\textendash 10\,\%}}.

وبضبط كمية ما قبل التخمير يمكنك التأثير في المدة التي يحتاجها عجينك في مرحلة التخمير الأول. فكلما زادت كمية البادئة التي تستخدمها تسارعت هذه العملية، وكلما قلّت تباطأت. فبكمية أكبر من البادئة تدخل كائنات دقيقة أكثر إلى عجينك الرئيسي، وكلما زادت هذه الكمية تسارع معدل التخمير في عجينك.

والعامل الآخر المؤثر في معدل التخمير هو درجة حرارة عجينك: فكلما ارتفعت درجة الحرارة تسارعت العملية، وكلما انخفضت تباطأت.

ما دام الغذاء متاحًا تتكاثر الكائنات الدقيقة وتزداد كمية. وهذه العملية محدودة بذاتها: فهي تتوقف حين لا يبقى غذاء متاحًا. ويمكن المقارنة هنا بصناعة النبيذ حيث تتجرع الخميرة في النهاية وتموت مع ارتفاع مستوى الإيثانول. فيخلق الإيثانول بيئة لا يمكن لكائنات دقيقة أخرى أن تنضم فيها إلى المأدبة. ويحدث الأمر نفسه مع الحموضة التي تنتجها البكتيريا: فالحموضة العالية تُبطئ عملية التخمير وتمنع كائنات دقيقة جديدة من دخول النظام.

في البداية تنتقل خصائص بادئتك إلى عجينك الرئيسي، ثم تتكيف الكائنات الدقيقة مع البيئة الجديدة مع الوقت. فإذا كانت بادئتك بكتيرية جدًّا فسيكون تخمير عجينك الرئيسي بكتيريًّا هو الآخر، وينتهي بك الأمر بعجين ليس منفتحًا كما يمكن أن يكون، ومذاقه حامض جدًّا لمعظم الناس، حامض أكثر من اللازم.

ولو استخدمت قيمة متطرفة نحو \LR{90\,\%} من البادئة بالنسبة إلى دقيقك، لما تبقى مجال واسع للكائنات الدقيقة كي تضبط توازنها في العجين الرئيسي. أما لو استخدمت \LR{1\,\%} فقط، أمكنت كائناتك الدقيقة أن تنمو من جديد داخل العجين إلى توازن مرغوب. وعليك أيضًا أن تأخذ في الاعتبار أن الكمية العالية من البادئة تعني تلقيحًا كبيرًا بدقيق مخمَّر من قبل. وكما ذكرنا سابقًا، تفكك الإنزيمات العجين؛ فكلما ارتفعت هذه القيمة زادت كمية الدقيق المخمَّر المتفكك. والتخمير الأطول من اللازم ينتج دائمًا عجينًا لزجًا جدًّا لا يمكن التعامل معه. فكلما زادت كمية البادئة وصلت إلى هذه النقطة أسرع. ولو استخدمت كمية ضئيلة جدًّا من البادئة فقد يتفكك دقيقك طبيعيًّا قبل أن يبلغ التخمير المرحلة المرجوة. ويمكن ملاحظة ذلك عند استخدام كمية صغيرة نحو \LR{1\,\%} من البادئة الحامضة: فكمية الكائنات الدقيقة المضافة الضئيلة لن تستطيع التكاثر بسرعة كافية قبل أن يفكك إنزيم البروتييز عجينك تفككًا تامًّا.

وكما شرحنا سابقًا، فإن مفتاح صنع خبز رائع هو تخمير بطيء لكن ليس بطيئًا أكثر من اللازم. فالإنزيمات تحتاج إلى وقت لتفكك عجينك. وبأخذ كل ذلك في الاعتبار، أحاول أن أستهدف زمن تخمير نحو \LR{8} إلى \LR{12} ساعة. وهذه تبدو النقطة الذهبية لمعظم الأدقائق التي عملت معها. ولتحقيق ذلك أستخدم نحو \LR{5\,\%} من البادئة الحامضة في الصيف (حيث درجات حرارة المطبخ حول \LR{25 °C (77 °F)}). أما في الشتاء فأختار نحو \LR{10\,\%} حتى \LR{20\,\%} من البادئة الحامضة (حيث حرارة المطبخ نحو \LR{20 °C (68 °F)}). ويتيح لي ذلك استخدام بادئة حامضة ليست في حالة مثالية. فكما شرحنا سابقًا، عجين خبزك في جوهره بادئة عملاقة، ومعدل التلقيح المنخفض يتيح للبادئة أن تنمو داخل عجينك الرئيسي إلى توازن مرغوب. وعلاوة على ذلك يتوفر للإنزيمات وقت كافٍ لتفكك الدقيق. ويتيح لي ذلك أيضًا الاستغناء تمامًا عن خطوة أوتولايز المزعومة (المزيد في القسم التالي)، وهذا يبسّط العملية كلها تبسيطًا كبيرًا.

\section{أوتولايز}%
\label{sec:autolysis}

يعني أوتولايز مزج الدقيق والماء فقط وترك الخليط راكدًا مدة تقارب \LR{30} دقيقة إلى عدة ساعات. وبعد اكتمال هذه العملية تُضاف البادئة الحامضة والملح إلى العجين.\footnote{جرّبتُ إضافة الملح في بداية عملية أوتولايز وفي نهايتها ولم ألحظ أي فرق. وبحسب فهمي الحالي، فإن أهمية إضافة الملح لاحقًا تبدو أسطورة.} فالمدة الكلية لتماس الدقيق والماء تتمدد، وبذلك تحصل على التفاعلات الإنزيمية النافعة التي تحسّن مذاق العجين وخصائصه. ولا أنصح بأوتولايز لأنه يضيف خطوة غير ضرورية إلى العملية، بل أنصح بتقنية الفيرمنتوليز \textenglish{(fermentolysis)} التي يتناولها القسم التالي من هذا الكتاب.

وآثار أوتولايز مثيرة للاهتمام جدًّا. جرّب أن تمزج الدقيق والماء فقط وتترك الخليط يومًا كاملًا. وأثناء اليوم افحص قوام عجينك، وحاول أن تمدّه. وإن جرأت يمكنك تذوّق العجين طوال اليوم. فمع كل ساعة يصبح عجينك أكثر قابلية للمد، ويسهل مدّه أكثر، وفي الوقت نفسه يبدأ مذاقه في التحول إلى الحلو حلوًا بعد حلو. فالإنزيمان البروتييز والأميلاز يؤديان عملهما. وتُستخدم العملية نفسها عند صنع حليب الشوفان: فبترك الخليط راكدًا بعض الوقت تعمل الإنزيمات على الشوفان، فيُستشعر المذاق أحلى وأكثر استحسانًا. وتتسارع هذه العملية أكثر كلما زادت نسبة الحبة الكاملة في دقيقك؛ فالقشرة تحتوي إنزيمات أكثر. وفي النهاية ستمزّق شبكة الغلوتين ويصطبغ عجينك، وهذا أعداؤك الأول في الخبز الحامض القمحي. فحين يحدث ذلك يصبح عجينك مسربًا ويفقد كل الغاز الثمين الذي أُنتج أثناء التخمير. وعليك أن تجد التوازن الصحيح بحيث يتفكك عجينك بما يكفي دون إفراط.

وعند استخدام معدل تلقيح عالٍ نحو \LR{20\,\%} من البادئة الحامضة يمكن أن يكون تخميرك سريعًا جدًّا. فعند \LR{25 °C} يمكن أن يكتمل في \LR{5} ساعات فقط. وإذا خمّرت مدة أطول أصبح عجينك مسربًا. وفي الوقت نفسه، خلال هذه الساعات الخمس، لم تتفكك الإنزيمات الدقيق بما يكفي، وهذا يعني أن العجين قد لا يكون مرنًا كما ينبغي. وعلاوة على ذلك لم تتحرر سكريات كافية، وبذلك لا يكون مذاق الخبز بعد الخَبز جيدًا بما يكفي.\footnote{لم أشاهد بعدُ دراسات تتناول سرعات الإنزيمات بحسب درجة الحرارة، لكني أفترض أنه كلما ارتفعت الحرارة تسارعت هذه التفاعلات، حتى حدٍّ معين تفكك عنده الإنزيمات بالحرارة.} لهذا يختار الخبازون أوتولايز. فيبدأ أوتولايز التفاعلات الإنزيمية قبل أن يبدأ تخمير الكائنات الدقيقة. وبهذه الطريقة، بعد ساعتين من أوتولايز (مثلًا) و\LR{5} ساعات من التخمير يكون العجين في الحالة المثالية قبل بدء التخمير النهائي.

وحين تحاول مزج الملح والبادئة في عجين الدقيق والماء ستلاحظ صعوبة الأمر؛ فتظن أنك مضطر إلى العجن من الصفر مرة أخرى، وستقضي وقتًا أطول في خلط العجين.

ولهذا السبب أدعو بحرارة إلى اعتماد نهج الفيرمنتوليز الذي يبسّط عملية الخلط والعجن تبسيطًا كبيرًا.

\section{الفيرمنتوليز}%
\label{sec:fermentolysis}

يخلق الفيرمنتوليز خصائص العجين النافعة نفسها التي يخلقها أوتولايز، لكن دون صداع مزج عجينك مرتين. تفعل ذلك بتمديد زمن تخمير عجينك: بدلًا من أوتولايز مدته ساعتين وتخمير أول مدته \LR{5} ساعات، تختار فترة تخمير كلية مدتها \LR{7} ساعات.

وللقيام بذلك تستخدم كمية أقل من البادئة الحامضة. فالوصفة التقليدية التي تتضمن خطوة أوتولايز قد تستدعي \LR{20\,\%} من البادئة الحامضة، وما عليك إلا خفض هذه القيمة إلى \LR{\textenglish{5\,\%\textendash 10\,\%}}. والخيار الآخر هو وضع العجين في بيئة أبرد، وبذلك تبطئ سرعة تكاثر كائناتك الدقيقة.

\begin{table}[htb]
    \centering
        % table-starter-usage-activity-ar.tex — Arabic edition of upstream
% tables/table-starter-usage-activity.tex (ch. 9, جدول 9.1)
%
% Original: The Sourdough Framework by Hendrik Kleinwächter,
% the-sourdough-framework.com, CC BY-SA 4.0. The Arabic edition inherits
% CC BY-SA 4.0.
%
% RTL mirror per the settled table mechanics (unit 9): the column order
% is reversed in the declared LTR tabular (upstream {@{}cSS@{}} becomes
% {@{}ccc@{}} with the starving column leftmost) so the °C/°F label
% column renders rightmost and the Amount (%) span covers the two
% value columns it covers upstream (declared cols 1--2 here, \cmidrule
% moved with them). Every non-empty cell wrapped {\textdir TRT ...}
% per the unit 9 cell-direction convention; numerals ride \LR{} islands
% (siunitx output must stay LTR inside RTL paragraphs); the °C/°F
% header and the "30 / 86" temperature pairs are Latin/symbol content
% and ride \LR{} islands whole. The recently-fed/starving column order
% and all cell values follow the accepted edition's render of p. 84
% (values 5/10/15 recently fed, 2.5/5/10 starving); the visual column
% order is render-verified against accepted p. 84 in the unit battery.

\begin{tabular}{@{}ccc@{}}
\toprule
\multicolumn{2}{c}{\thead{{\textdir TRT الكمية \LR{(\%)} للبادئة}}} & \\
\cmidrule(lr){1-2}
\thead{{\textdir TRT جائعة}} & \thead{{\textdir TRT تغذيت حديثًا}} & \thead{{\textdir TRT \LR{°C / °F}}} \\
\midrule
{\textdir TRT \LR{\num{2.5}}} & {\textdir TRT \LR{\num{5}}}  & {\textdir TRT \LR{30 / 86}} \\
{\textdir TRT \LR{\num{5}}}   & {\textdir TRT \LR{\num{10}}} & {\textdir TRT \LR{25 / 77}} \\
{\textdir TRT \LR{\num{10}}}  & {\textdir TRT \LR{\num{15}}} & {\textdir TRT \LR{20 / 68}} \\
\bottomrule
\end{tabular}

        \caption[كمية البادئة الحامضة]{جدول يوضح كمية البادئة الحامضة التي ينبغي استخدامها بحسب درجة الحرارة ومستوى نشاط البادئة.}
\end{table}

وبحسب تجربتي وبحسب بادئتي، يأخذ خبزي المثالي دائمًا نحو \LR{8} إلى \LR{12} ساعة أثناء التخمير الأول. وبحسب الأوقات المتاحة لي في اليوم أستخدم كمية أكبر أو أصغر من البادئة. فإذا أردت اكتمال التخمير في \LR{8} ساعات اخترت بادئة حامضة بنسبة \LR{10\,\%}، وإذا أردته جاهزًا في \LR{12} ساعة اخترت بادئة أقل نحو \LR{5\,\%}. ما عليك إلا مزج جميع المكونات ويبدأ تخميرك: فالإنزيمات والكائنات الدقيقة يبدآن عملهما. وفي يوم صيفي حار جدًّا لم تعد الكميات المذكورة تفيد: فببادئة \LR{10\,\%} يصبح العجين نفسه جاهزًا في \LR{5} ساعات حتى نقطة اللاعودة، وكل ساعة إضافية تسبب تفكك العجين تفككًا زائدًا. وفي هذه الحالة سأختار بادئة حامضة \LR{5\,\%} لإبطاء العملية كلها والعودة إلى نافذة \LR{8} إلى \LR{12} ساعة. وإذا كان الجو حارًّا جدًّا فقد أستخدم \LR{1\,\%} فقط من البادئة الحامضة\footnote{خذ هذه القيم بحبة ملح، فهي تعتمد على دقيقك وبادئتك الحامضة، وهي قيم تحتاج إلى التجربة فيها بنفسك. وبعد خَبز عدد من الأرغفة ستتمكن من قراءة عجينك تحسنًا كبيرًا.}. وعليك أن تلعب بالأوقات بنفسك. لكن بدلًا من الاعتماد على التوقيت سأريك نهجًا أفضل وأدق بكثير باستخدام عينة تخمير، وسنتناول ذلك لاحقًا في هذا الفصل.

وحتى في عجائن الخميرة لم أعد أستخدم أوتولايز؛ أكتفي بخفض كمية الخميرة التي أستخدمها. واعتماد الفيرمنتوليز يوفر عليك الوقت ويبسّط عملية صنع الخبز. وكما ذكرنا في الفصول السابقة، سرّ صنع الخبز الرائع هو تخمير بطيء لكن ليس بطيئًا أكثر من اللازم.

\section{قوة العجين}

قوة العجين هي تسمية بذخ لوصف عملية عجن الخبز. فكلما انتظرت وعجنّت تقوّت روابط الغلوتين في عجينك، وصار العجين أكثر مرونة وتماسكًا. وهذا أساس حصر الغازات كلها أثناء عملية التخمير؛ فبدون شبكة الغلوتين سيتسرب الغاز من عجينك منتشرًا.

\begin{flowchart}[htb]
% fig-kneading-process-ar.tex — Arabic edition of upstream
% figures/fig-kneading-process.tex (ch. 9, المخطط الانسيابي 9.3)
%
% Original: The Sourdough Framework by Hendrik Kleinwächter,
% the-sourdough-framework.com, CC BY-SA 4.0. The Arabic edition inherits
% CC BY-SA 4.0.
%
% R1-FIGURES route: re-derived from the upstream source
% figures/fig-kneading-process.tex (node styles from
% figures/flowcharts_tikz.tex, node/edge geometry verbatim), with the
% label text recovered from the accepted edition's render of p. 86
% (book-ar-fixed.pdf, b78c2a28…). Line splits follow the accepted
% render (اخلط/المكونات pattern: امزج/مكونات الوصفة، انتظر/15 دقيقة،
% اعجن/5 دقائق، اختبار/النافذة؟، الباسيناج/لمزيد من الماء؟، اصنع كرة/
% عجين ملساء، جرّبة/الأليكوت؟، ابدأ/التخمير الأول). Arabic label text
% wrapped in \textarabic{}; digit runs ride \LR{} islands (Ruling 1
% uniform-Pagella figure face, ACCEPTED 2026-11-15). Decision edge
% labels نعم/لا follow the accepted render.
%
% ACCEPTED-EDITION DIVERGENCE (2026-11-16): the accepted edition renders
% the second wait node (upstream "Wait 15 minutes", wait2) as an EMPTY
% box — its text is absent from both the accepted text layer and the
% render of acc. p. 86. The node is kept (box drawn) with empty text to
% match the accepted edition; upstream geometry retained. The accepted
% render also displaces the knead node's «5 دقائق» line outside its box
% (accepted-build artifact); here it stays inside the اعجن node, letters
% unchanged.

\begin{tikzpicture}[node distance = 4cm, auto]
  \node [start] (init) {\textarabic{امزج}\\ \textarabic{مكونات الوصفة}};
  \node [block, right of=init] (wait1) {\textarabic{انتظر}\\ \LR{15} \textarabic{دقيقة}};
  \path [line] (init) -- (wait1);
  \node [block, right of=wait1] (knead) {\textarabic{اعجن}\\ \LR{5} \textarabic{دقائق}};
  \path [line] (wait1) -- (knead);
  % accepted edition renders this node's box empty (text layer + render of acc. p. 86):
  \node [block, right of=knead] (wait2) {};
  \path [line] (knead) -- (wait2);
  \node [decision, below of=wait2, node distance=4cm] (windowpane_test) {\textarabic{اختبار}\\ \textarabic{النافذة؟}};
  \path [line] (wait2) -- (windowpane_test);
  \path [line] (windowpane_test.east) -- node{\textarabic{لا}} ++(1, 0) -- node{} ++(0, 5.7) -|  (knead.north);
  \node [decision] at (windowpane_test -| knead) (more_water) {\textarabic{الباسيناج}\\ \textarabic{لمزيد من الماء؟}};
  \path [line] (windowpane_test) -- node{\textarabic{نعم}} (more_water);
  \node [block] at (wait1 |- more_water) (add_water) {\textarabic{أضف الماء}};
  \path [line] (more_water) -- node{\textarabic{نعم}} (add_water);
  \path [line] (add_water.north east) -- (knead.south west);
  \node [decision, below of=more_water, node distance=4.5cm] (dough_sample) {\textarabic{جرّبة}\\ \textarabic{الأليكوت؟}};
  \path [line] (more_water) -- node{\textarabic{لا}} (dough_sample);
  \node [block] at (dough_sample -| wait2) (dough_ball) {\textarabic{اصنع كرة}\\ \textarabic{عجين ملساء}};
  \path [line] (dough_sample) -- node{\textarabic{لا}} (dough_ball);
  \node [block, below of=dough_sample] (extract_sample) {\textarabic{استخرج جرّبة}};
  \path [line] (dough_sample) -- node{\textarabic{نعم}} (extract_sample);
  \path [line] (extract_sample.north east) -- (dough_ball.south west);
  \node [success, below of=dough_ball] (begin_bulk) {\textarabic{ابدأ}\\ \textarabic{التخمير الأول}};
  \path [line] (dough_ball) -- (begin_bulk);
\end{tikzpicture}

\caption[عملية تطوير الغلوتين في عجين قمحي]{عملية تطوير الغلوتين في عجين قمحي.}
\label{fig:wheat-sourdough-kneading-process}
\end{flowchart}

قد يبدو الأمر غريبًا، لكن أهم جزء في العجن هو الانتظار. فبالانتظار تتيح لدقيقك امتصاص الماء، وبهذا تتشكل روابط الغلوتين في عجينك تلقائيًا ويصبح العجين أكثر مرونة. فلعله تعجن \LR{10} دقائق في البداية لتتفاجأ أن العجن \LR{5} دقائق مع انتظار \LR{15} دقيقة له المفعول نفسه.

فبروتينا الغلوتين الغلوتينين والغليادين يرتبطان فورًا تقريبًا بعد ترطبهما. وتتيح روابط ثنائي الكبريت للأجزاء الطويلة من الغلوتينين أن تنضم بعضها إلى بعض مشكّلة جزيئات متينة قابلة للمد. فيضيف الغلوتينين القوة، بينما تتيح جزيئات الغليادين الأكثر إحكامًا للعجين أن يجري مثل السائل. وفي النهاية، كلما انتظرت أكثر تحوّلت شبكة الغلوتين أكثر إلى بنية شبه عنكبوتية، وهي ما يحصر الغازات أثناء عملية التخمير~\cite{how+does+gluten+work}.

\begin{figure}[htb]
  \centering
  \includegraphics[width=\textwidth]{dough-strength-no-kneading}
  \caption[قوة العجين مع الوقت دون عجن]{تصور تخطيطي للتطور التلقائي للغلوتين. فالعجائن هنا غير معجنّة، ولم تُخلط فيها إلا في البداية. لاحظ كيف تتدهور قوة العجين مع الوقت بينما تفكك الإنزيمات الدقيق. ويتسارع هذا التأثير في الخبز الحامض بسبب التحلل البروتيني للغلوتين الذي تسببه البكتيريا.}
  \label{fig:wheat-yeast-sourdough-degradation}
\end{figure}

عملية النقع ينبغي أن تتمدد كلما زاد استخدام دقيق القمح كامل الحبة. فغرض قشرة حبة القمح الخارجية أن تمتص الماء بأسرع ما يمكن. فتُنشَّط الإنزيمات وتبدأ عملية الإنبات. ولهذا يقل الماء المتاح لتطوّر روابط الغلوتين. فإما أن تنتظر مدة أطول قليلًا أو تمضي وتستخدم ماءً أكثر للعجين قليلًا.

وهذا هو المبدأ نفسه الذي تتبعه وصفات \textenglish{no-knead} الشهيرة: فبصنع عجين أقل ترطيبًا والانتظار تتشكل شبكة غلوتينك تلقائيًا. وما تزال مضطرًا إلى خلط المكونات وتجنيسها. وتنتظر بضع دقائق فقط لترى عجينك قد طوّر قوة عجين مذهلة دون أي عجن إضافي.\footnote{جرّب ذلك بنفسك؛ فالتشكل التلقائي لشبكات الغلوتين ظاهرة مدهشة ما تزال تأسرني في كل مرة أصنع فيها عجينًا.}

فإذا رطبت عجينك رطوبة زائدة في البداية صار من الأصعب على سلاسل الغلوتين أن تتشكل. فالجزيئات ليست متقاربة بعضها من بعض في عجين رطب مقارنة بعجين أكثر صلابة، ويزوم عليها أن تنسجم وتشكل البنية الشبكية. ولهذا يظل الأسهل دائمًا البدء بترطيب أقل ثم زيادة كمية الماء عند الحاجة. وهذا معروف عمومًا بطريقة الباسيناج \textenglish{(Bassinage)}. فروابط الغلوتين تكون قد تشكلت عند الترطيب المنخفض، ثم يمكن جعلها أكثر قابلية للمد بإضافة الماء والعجن من جديد. وهذه حيلة رائعة لصنع عجين أكثر قابلية للمد بدقيق أقل غلوتينًا~\cite{bassinage+technique}.

وعند العجن بالآلة اعتمد التقنية نفسها المعروضة في المخطط~\ref{fig:wheat-sourdough-kneading-process}. فابدأ بسرعة منخفضة؛ فهذا يساعد عملية التجنيس. وبعد الانتظار ليتحصل الدقيق على الماء انتقل إلى إعداد سرعة أعلى. ومن العلامات الجيدة على شبكة غلوتين مطورة جيدًا أن يفلت عجينك من الوعاء؛ فمرونة الغلوتين تفوق رغبة العجين في الالتصاق بالوعاء.

\begin{figure}[htb]
  \centering
  \includegraphics[width=\textwidth]{dough-strength-sourdough}
  \caption[قوة العجين مع الوقت مع العجن]{تصور تخطيطي لتطور الغلوتين في الخبائز الحامضة بتقنيات عجن مختلفة. ويمكن استخدام مزيج من التقنيات لبلوغ أقصى قوة عجين.}
  \label{fig:dough-strength-sourdough}
\end{figure}

وعمومًا، كلما خلقتم قوة عجين أكبر شعرتم بأن عجينكم أقل لزوجة. فبما أن العجين يصبح متماسكًا، لن يلتصق بيدك كثيرًا. وهذه مشكلة شائعة يواجهها المبتدئون: فالعجين اللاصق كثيرًا ما يكون علامة على شبكة غلوتين غير مطورة بما يكفي.

\begin{figure}[htb]
  \centering
  \includegraphics[width=\textwidth]{dough-surface-touchpoints}
  \caption[لمس سطح العجين]{تصور تخطيطي لكيفية أن سطح العجين الخشن يخلق نقاط تماس أكثر من سطح العجين الأملس. فعند لمس السطح الخشن يتورم العجين ويتلامس مع مساحات أكثر من يدك.}
  \label{fig:dough-touch-points}
\end{figure}

العجن أكثر نافع عمومًا في كل الحالات تقريبًا لأنه ينتج شبكة غلوتين أقوى. غير أنك عند صنع خبز الحليب الطري قد تفضل عجينًا أكثر قابلية للمد من البداية. وفي هذا السيناريو قد يؤدي العجن المفرط إلى خبز نهائي أكثر مطاطية، وهو أمر غير مرغوب إذا كنت تستهدف قوامًا أكثر فتوحة. ويمكن تحقيق هذه الفتوحة بالعجن الأقل. ورغم أن هذا استثناء، فإن العجن السليم لعجائنك القمحية نصيحة عامة.

وعند استخدام العجنة الكهربائية قد تصادف مشكلة العجن الزائد، لكن هذا شبه مستحيل عمليًّا. فحتى بعد العجن \LR{30} دقيقة على سرعة متوسطة، نادرًا ما كانت عجائني معجنّة أكثر من اللازم. وفي اللحظة التي تعجن فيها أكثر من اللازم يمكن أن يبدأ لون العجين في التغير، لكنك تلاحظ ذلك غالبًا أثناء الخَبز: فالرغيف الناتج يبدو باهتًا وأبيض. والسبب أن عجن العجين يسبب الأكسدة الضرورية لتطوير الغلوتين؛ أما إذا عُجن العجين كثيرًا جدًّا فقد تتلف المركبات التي تسهم في نكهة الخبز ورائحته ولونه، بما يؤثر سلبًا في جودة الخبز~\cite{oxidization+dough}.

والخطوة الأخيرة قبل بدء التخمير الأول هي صنع كرة عجين ملساء. فبضمان أن سطح عجينك أملس ستكون نقاط التماس أقل عند لمس العجين. انظر الشكل~\ref{fig:dough-touch-points} لتصور تخطيطي لكيفية لمس يدك للعجين الخشن والملساء.

وبالسطح الملساء يلتصق عجينك بيديك أقل، ويصبح تطبيق الطيّات والشدّات لاحقًا أسهل بكثير. وبدون سطح ملساء يصبح العجين شبه غير قابل للعمل، ويصبح طي العجين لاحقًا مهمة مستحيلة. وهذا خطأ متكرر أراه يرتكبه خبازون مبتدئون كثيرون.

\begin{figure}[htb]
  \centering
  \includegraphics[width=0.8\textwidth]{dough-ball-steps}
  \caption[صنع سطح أملس]{تحويل كتلة عجين لاصقة إلى عجين بسطح أملس. الغرض هو تقليص نقاط تماس السطح بيديك لتصبح لزوجة العجين أقل أثناء عمله.}
  \label{fig:dough-ball-steps}
\end{figure}

ولجعل سطح العجين ملساء، ضع عجينك على لوح خشبي أو على رخامة مطبخك، واسحب العجين براحة يدك على السطح، ويمكنك الاستعانة بكاشط العجين هنا. اسحب العجين نحوك مع الحرص على بقاء مركز العجين العلوي في مكانه. ويساعد أن تضع يدك الثانية بلطف فوق العجين بحيث تتحرك كتلة العجين مع الحفاظ على اتجاهها. وحين يقترب العجين كله كثيرًا من حافة الوعاء أو الرخامة أعده بلطف بيديك الاثنتين. وبفعل ذلك تمدّ طبقة الغلوتين المحيطة الخارجية. ولهذا من المهم ألا تستخدم أي دقيق أثناء هذه العملية؛ فباستخدام الدقيق لن تتمكن من سحب العجين على السطح، وبذلك لا تستطيع مدّ الغلوتين. وتخيل دائمًا أنك تلمس شيئًا لاصقًا تمامًا، فبذلك ستحاول تلقائيًّا لمس العجين بأقل قدر ممكن. وواصل تكرار العملية حتى ترى أن للعجين سطحًا جميلًا ملساء. وينبغي أن يبدو العجين النهائي مثل العجين المعروض في الشكل~\ref{fig:dough-ball-steps}.

فإذا تمزقت طبقة الغلوتين الخارجية فقد مددت عجينك أكثر من اللازم. وفي هذه الحالة خذ استراحة \LR{10} دقائق مع ترك عجينك على رخامة المطبخ، فيتاح للغلوتين أن يعيد الارتباط ويلتئم. وكرر العملية نفسها وينبغي أن تختفي المناطق الخشنة المتضررة.

وتُستخدم تقنية تدوير العجين نفسها لاحقًا أثناء عملية التشكيل المسبق. وبعد بناء قوة العجين يصبح لديك كل الوقت الذي تحتاجه لممارسة التدوير. دوّر العجين قدر الإمكان حتى يتمزق، ثم انتظر \LR{10} دقائق المذكورة سابقًا وكرر. أما لاحقًا فلن يكون لديك مجال للخطأ، ولا بد أن تكون تقنيتك في القمة؛ فالعجين المشكَّل مسبقًا أكثر من اللازم قد لا يتعافى.

\section{التخمير الأول}
\label{sec:bulk-fermentation}

بعد مزج البادئة في عجينك تبدأ المرحلة التالية من العملية المعروفة بالتخمير الأول. ويُستخدم مصطلح الكتلة لأن المخبزات تخمّر أرغفة متعددة معًا في كتل. أما إذا كنت خبازًا منزليًّا فقد تخمّر رغيفًا واحدًا. وينتهي التخمير الأول حين تقسم وتشكّل مسبقًا، أو تشكّل أرغفتك النهائية أو رغيفك النهائي مباشرةً.

وأصعب جزء في صنع الخبز الحامض هو ضبط عملية التخمير. فالتخمير الأول مدة كافية وليست زائدة هو العامل الحاسم في صنع خبز رائع في البيت. فحتى بتقنيات تشكيل وخَبز سيئة ستتمكن من صنع خبز ممتاز بمجرد إتقان عملية التخمير الأول وحدها.

وبتخمير أول قصير جدًّا ستُستشعر لُبّك لزجةً، وسيضم لُبّك جيوبًا كبيرة من الهواء يُشار إليها عمومًا باسم الحفر. أما التخمير الأطول من اللازم فينتج عنه تفكك العجين تفككًا زائدًا: فيلتصق العجين الناتج بسلة التخمير ويتمدد أثناء الخَبز إلى بنية تشبه الفطيرة.

والسر هو إيجاد النقطة الذهبية بين تخمير أول ليس قليلًا وليس كثيرًا. وأنصح دائمًا بدفع العجين أكثر نحو التخمير الأطول؛ فنكهة الخبز الناتج أفضل مقارنة بالعجين الشاحب الناقص التخمير.

\begin{table}[htb]
    \centering
        % table-fermentation-effects-ar.tex — Arabic edition of upstream
% tables/table-fermentation-effects.tex (ch. 9, جدول 9.2)
%
% Original: The Sourdough Framework by Hendrik Kleinwächter,
% the-sourdough-framework.com, CC BY-SA 4.0. The Arabic edition inherits
% CC BY-SA 4.0.
%
% RTL mirror per the settled table mechanics (unit 9): the column order
% is reversed in the declared LTR tabular (upstream
% {@{}>{\bfseries}p{0.12\textwidth}p{0.273\textwidth}
% p{0.273\textwidth}p{0.273\textwidth}@{}} becomes
% {@{}p{0.273\textwidth}p{0.273\textwidth}p{0.273\textwidth}
% >{\bfseries}p{0.12\textwidth}@{}}) so the row-label column renders
% rightmost and keeps its bold \bfseries prefix, and the Fermentation
% span covers the three value columns it covers upstream (declared
% cols 1--3, \cmidrule moved with them; the perfect column sits
% leftmost, the too-short column innermost next to the labels). Every
% non-empty cell wrapped {\textdir TRT ...} per the unit 9
% cell-direction convention; p-column cells are RTL paragraphs, so
% multi-sentence cells read naturally. The guillemet word «فوهات»
% (upstream ``craters'') keeps the accepted edition's guillemets.
% Cell text follows the accepted edition's render of p. 91 line for
% line; the caption lives in the chapter file.

\begin{tabular}{@{}p{0.273\textwidth}p{0.273\textwidth}p{0.273\textwidth}>{\bfseries}p{0.12\textwidth}@{}}
\toprule
\multicolumn{3}{c}{\thead{{\textdir TRT التخمير}}} & \\
\cmidrule(lr){1-3}
\thead{{\textdir TRT مثالية}} & \thead{{\textdir TRT طويلة جدًا}} & \thead{{\textdir TRT قصيرة جدًا}} & \\
\midrule
{\textdir TRT لُبّ مخبوز بانتساوية؛ قد يُدرَك رطبًا لكنه ليس لزجًا.} &
{\textdir TRT قد يُدرَك اللُّبّ لزجًا لأن معظم الغلوتين قد تفكك.} &
{\textdir TRT مناطق غير مخبوزة لزجة باتجاه أسفل الخبز.} &
{\textdir TRT القوام اللُّبّ} \\
{\textdir TRT فجوات موزعة بانتساوية بلا «فوهات».} &
{\textdir TRT فجوات صغيرة كثيرة موزعة بانتساوية.} &
{\textdir TRT فجوات ضخمة على نحو مبالغ فيه داخل اللُّبّ على شكل «فوهات».} &
{\textdir TRT الفجوات الهوائية} \\
{\textdir TRT طابع نكهة متوازن، ليس باهتًا أكثر مما ينبغي ولا حامضًا أكثر مما ينبغي؛ وحسب البادئة لمسات خلّ أو لاكتيك.} &
{\textdir TRT طابع حمضي قوي؛ تطغى الحمضية عند التذوق.} &
{\textdir TRT مذاق باهت محايد.} &
{\textdir TRT المذاق} \\
{\textdir TRT مزيج رائع من الأقوام.} &
{\textdir TRT تماسك جيد، لكن اللُّبّ ليس منفوشًا كما يمكن أن يكون.} &
{\textdir TRT قوام رديء إجمالًا.} &
{\textdir TRT القوام} \\
{\textdir TRT انتفاخ عمودي رائع؛ ينمو العجين صعودًا أكثر من نموه جانبًا.} &
{\textdir TRT بنية مسطحة جدًا تشبه الفطيرة بعد الخَبز.} &
{\textdir TRT انتفاخ عمودي، سببه غالبًا تبخر الماء وانتفاخ العجين.} &
{\textdir TRT الانتفاخ في الفرن} \\
\bottomrule
\end{tabular}
%
        \caption[مراحل التخمير الحامض]{المراحل المختلفة للتخمير الحامض وآثارها في اللُّبّ والفتحات والقوام والمذاق الكلي.}
\end{table}

وأقبح ما يمكنك فعله أثناء تخمير الخبز الحامض أن تعتمد على اقتراحات الوصفة للتوقيت. ففي \LR{99\,\%} من الحالات لن يناسبك التوقيت؛ فكاتب الوصفة غالبًا له دقيق مختلف وبادئة حامضة مختلفة بمستويات نشاط مختلفة. وعلاوة على ذلك قد تكون درجة حرارة بيئة التخمير مختلفة. والتغييرات الصغيرة في متغير واحد تنتج جدول توقيت مختلف تمامًا: ففرق ساعة أو ساعتين ينتج عجينًا لم يخمّر مدة كافية، أو يتحول إلى فطيرة مخمّرة عملاقة لاصقة. وهذا أحد أسباب تفضيل صناعة الخَبز الحالية عجائن الخميرة فقط: فبإزالة البكتيريا من التخمير تصبح العملية كلها أكثر تنبؤًا بكثير، ويصبح مجال الخطأ (كما في الشكل~\ref{fig:wheat-yeast-sourdough-degradation}) أوسع بكثير، وتصبح العجائن مثالية للصناعة في الآلات.

\begin{flowchart}[htb]
  \centering
  % fig-bulk-fermentation-ar.tex — Arabic edition of upstream
% figures/fig-bulk-fermentation.tex (ch. 9, المخطط الانسيابي 9.4)
%
% Original: The Sourdough Framework by Hendrik Kleinwächter,
% the-sourdough-framework.com, CC BY-SA 4.0. The Arabic edition inherits
% CC BY-SA 4.0.
%
% R1-FIGURES route: re-derived from the upstream source
% figures/fig-bulk-fermentation.tex (node styles from
% figures/flowcharts_tikz.tex, node/edge geometry verbatim including the
% no-double-line draw tricks), with the label text recovered from the
% accepted edition's render of p. 92 (book-ar-fixed.pdf, b78c2a28…).
% Line splits follow the accepted render (التخمير/الأول، افحص/العجين،
% افحص/زيادة الحجم، افحص/قيمة pH، افحص/رائحة العجين، هل العجين/جاهز؟،
% قسّم/وشكّل مسبقًا، هل سُطِّح/العجين؟، انتظر/60 دقيقة، الطيّات/
% والشدّات). Arabic label text wrapped in \textarabic{}; the pH symbol
% and the 60-minute digit ride \LR{} islands (Ruling 1 uniform-Pagella
% figure face, ACCEPTED 2026-11-15). The أو edge labels and نعم/لا
% decision labels follow the accepted render. The decision node
% هل سُطِّح العجين؟ carries a kasra+shadda on the س in the accepted
% render's text layer (هل س ّطح) — flagged for render verification in
% the unit battery.

\begin{tikzpicture}[node distance = 3cm, auto]
  \node [start] (init) {\textarabic{التخمير}\\ \textarabic{الأول}};
  \node [block, right of=init, node distance=4cm] (check_dough) {\textarabic{افحص}\\ \textarabic{العجين}};
  \node [block, right of=check_dough, node distance=4cm] (size_increase) {\textarabic{افحص}\\ \textarabic{زيادة الحجم}};
  \node [block, below of=size_increase, node distance=2cm] (ph_value) {\textarabic{افحص}\\ \textarabic{قيمة \LR{pH}}};
  \node [block, below of=ph_value, node distance=2cm] (smell) {\textarabic{افحص}\\ \textarabic{رائحة العجين}};
  \node [decision, right of=size_increase, node distance=4cm] (dough_ready) {\textarabic{هل العجين}\\ \textarabic{جاهز؟}};
  \node [success] at(dough_ready |- smell) (divide_preshape) {\textarabic{قسّم}\\ \textarabic{وشكّل مسبقًا}};
  \node [decision, above of=size_increase] (dough_flattened) {\textarabic{هل سُطِّح}\\ \textarabic{العجين؟}};
  \node [block, above of=check_dough] (wait_60_minutes) {\textarabic{انتظر}\\ \LR{60} \textarabic{دقيقة}};
  \node [block, above of=wait_60_minutes] (stretch_fold) {\textarabic{الطيّات}\\ \textarabic{والشدّات}};

  \path [line] (init) -- (check_dough);
  \path [line] (check_dough) -- (size_increase);
  % Tricks not to get double lines
  \path [line] (check_dough) ++(2, -2) -- node{\textarabic{أو}} (ph_value);
  \path [line] (check_dough) ++(2, 0) -- node{} ++(0, -4) -- node{\textarabic{أو}} (smell);
  \path [line] (check_dough) ++(2, -4) -- node{\textarabic{أو}} (smell);
  \path [line] (size_increase) -- (dough_ready);
  % Same tricks not to get double lines and also we do _not_ want arrows
  \path [draw, thick] (ph_value) -- node{} ++(2, 0);
  \path [draw, thick] (smell) -| node{} ++(2, 4);
  \path [line] (dough_ready) -- node{\textarabic{نعم}} (divide_preshape);
  \path [line] (dough_ready) |- node[right=3pt]{\textarabic{لا}} (dough_flattened);
  \path [line] (dough_flattened) |- node[right=3pt]{\textarabic{نعم}} (stretch_fold);
  \path [line] (dough_flattened) -- node{\textarabic{لا}} (wait_60_minutes);
  \path [line] (stretch_fold) -- (wait_60_minutes);
  \path [line] (wait_60_minutes) -- (check_dough);
\end{tikzpicture}

  \caption[عملية فحص التخمير الأول]{أثناء التخمير الأول تُخمَّر عجائن متعددة معًا في كتلة. ومن الجوانب الصعبة في صنع الخبز الحامض المنزلي تحديد متى تكتمل مرحلة التخمير هذه. ويوضح هذا المخطط خيارات متعددة متاحة لفحص سير التخمير الأول.}
  \label{fig:bulk-fermentation}
\end{flowchart}

سيقول لك الخبازون المتمرسون أن تحكم بحسب مظهر العجين وإحساسه. وهذا ينجح إذا كنت قد صنعت مئات الأرغفة، لكنه ليس خيارًا للخباز غير المتمرس. فكلما صنعت عجينًا أكثر ستتمكن من الحكم على حالة العجين باللمس.

وطريقتي المفضلة للمبتدئين هي استخدام جرّة الأليكوت \textenglish{(aliquot jar)}. فالأليكوت عينة تستخرجها من عجينك، تُستخرج بعد بناء قوة العجين الأولية. وتراقب زيادة حجم الأليكوت لتحكم في مستوى تخمير عجينك الرئيسي. فحين يتخمر عجينك يتخمر محتوى جرّة الأليكوت هو الآخر. وفي اللحظة التي تبلغ عينتك حجمًا معينًا يصبح عجينك الرئيسي جاهزًا للتشكيل والتخمير النهائي. وزيادة الحجم التي ينبغي أن تستهدفها تعتمد على الدقيق المتوفر لديك: فالدقيق الأعلى غلوتينًا يمكن تخميره مدة أطول. وعلمًا بأن نحو \LR{80\,\%} من بروتين دقيق القمح لديك هو غلوتين، فافحص عبوة دقيقك لترى نسبة البروتين. وقيمة زيادة الحجم الفعلية شديدة التغير بحسب تركيبة دقيقك. وأنصح بالبدء بزيادة حجم \LR{25\,\%} والتجربة حتى \LR{100\,\%} في الخبزات التالية، ثم تحديد القيمة التي ترضاك.

\begin{table}[htb]
    \centering
        % table-dough-size-increase-ar.tex — Arabic edition of upstream
% tables/table-dough-size-increase.tex (ch. 9, جدول 9.3)
%
% Original: The Sourdough Framework by Hendrik Kleinwächter,
% the-sourdough-framework.com, CC BY-SA 4.0. The Arabic edition inherits
% CC BY-SA 4.0.
%
% RTL mirror per the settled table mechanics (unit 9): the column order
% is reversed in the declared LTR tabular (upstream {@{}cc@{}} becomes
% {@{}cc@{}} with the size-increase column leftmost) so the flour
% protein content label column renders rightmost. Every non-empty cell
% wrapped {\textdir TRT ...} per the unit 9 cell-direction convention.
% The percentage ranges are literal-text ranges per the unit 9
% dash-font convention: \LR{\textenglish{N\textendash M\%}} keeps the
% en-dash and the percent sign in the Latin island; the "> 15%" cells
% ride \LR{} islands whole. Cell text follows the accepted edition's
% render of p. 93; the caption lives in the chapter file.

\begin{tabular}{@{}cc@{}}
\toprule
\thead{{\textdir TRT الزيادة النسبية في حجم العينة}} & \thead{{\textdir TRT نسبة البروتين في الدقيق}} \\
\midrule
{\textdir TRT \LR{25\%}}  & {\textdir TRT \LR{\textenglish{8\textendash 10\%}}} \\
{\textdir TRT \LR{50\%}}  & {\textdir TRT \LR{\textenglish{10\textendash 12\%}}} \\
{\textdir TRT \LR{100\%}} & {\textdir TRT \LR{\textenglish{12\textendash 15\%}}} \\
{\textdir TRT \LR{> 100\%}} & {\textdir TRT \LR{\textenglish{> 15\%}}} \\
\bottomrule
\end{tabular}
%
        \caption[زيادة الحجم مقابل محتوى البروتين]{قيم مرجعية لمقدار زيادة الحجم التي تستهدفها بجرّة الأليكوت بحسب محتوى البروتين في العجين.}
\end{table}

وجمال الأليكوت أنه بغض النظر عن درجة الحرارة المحيطة ستعرف دائمًا متى يصبح عجينك جاهزًا. فقد يكون العجين جاهزًا في \LR{8} ساعات في الصيف، لكنه قد يبلغ \LR{12} ساعة في الشتاء بسهولة. وستخمّر عجينك دائمًا في التوقيت الصحيح تمامًا.

\begin{figure}[htb]
  \centering
  \includegraphics[width=\textwidth]{aliquot-before-after}
  \caption[جرّة الأليكوت]{جرّة أليكوت لمراقبة سير تخمير العجين. استغرق الأمر \LR{10} ساعات حتى بلغ العجين زيادة حجم \LR{50\,\%}.}
\end{figure}

ورغم أن عينة الأليكوت مكنتني من خَبز أرغفة رائعة باستمرار، فهناك قيود يجب مراعاتها. من الجوهري استخدام وعاء أسطواني الشكل للحكم سليمًا على زيادة حجم العجين. وعلاوة على ذلك من الأساسي استخدام ماء بدرجة حرارة الغرفة عند صنع عجينك. فإذا كان الماء أدفأ فقد أبردت عينتك، بسبب حجمها الأصغر، أسرع من عجينك، وستتخمر أبطأ من عجينك الرئيسي. كذلك إذا استخدمت ماء باردًا أكثر من اللازم فإن عينتك ستدفأ أسرع من كتلة العجين الكبيرة، وفي هذه الحالة ستسبق عينتك عجينك الرئيسي، ومن المرجح أن توقف التخمير مبكرًا أكثر من اللازم. احرص على إبقاء العجين والأليكوت قريبين بعضهما من بعض، بل يضع بعض الناس الأليكوت في الوعاء نفسه، فيضمن ذلك أن كليهما في درجة حرارة البيئة نفسها. كما أن الأليكوت أقل موثوقية إذا تغيرت درجة الحرارة المحيطة تغيرًا كبيرًا خلال اليوم؛ ففي هذه الحالة ستتكيف عينتك أسرع من عجينك الرئيسي وستبقى القراءات منحرفة قليلًا. وإذا كنت تصنع كتلة عجين كبيرة بأكثر من \LR{10 kg} من الدقيق فالجرّة أقل موثوقية كذلك؛ فالتفاعلات البيوكيميائية التي تجري داخل عجينك ترفع حرارته، والتخمير نفسه طارد للحرارة، أي أنه ينتج حرارة.

والخيار الآخر الأغلى هو استخدام مقياس للأس الهيدروجيني \textenglish{(pH)} لمراقبة حالة تخمير عجينك. فبينما تتخمر بكتيريا حمض اللبنيك وحمض الخليك تتراكم الحموضة داخل عجينك. ويمكن قياس قيمة الحموضة (الأس الهيدروجيني) بمثل هذا المقياس: فكلما زادت الحموضة انخفضت قيمة الأس الهيدروجيني في عجينك. وسلم الأس الهيدروجيني لوغاريتمي، أي أن كل تغير وحدة يعني زيادة \LR{10} أضعاف في الحموضة. فقد يبدأ عجين الخبز الحامض التخمير عند \LR{pH 6.0} ثم يقارب \LR{pH 4.0} قبل الخَبز بمدة قصيرة، أي أن العجين نفسه أشد حموضة \LR{10} أضعاف من \LR{10} أضعاف (= \LR{100} ضعف) من بدايته. وباستخدام المقياس يمكنك دائمًا الحكم على حالة حموضة عجينك والتصرف تبعًا لذلك.

لاستخدام مقياس الأس الهيدروجيني بنجاح يلزمك إيجاد قيم تناسب عجينك. فبحسب بادئتك وماءك وتركيبة دقيقك تختلف قيم الأس الهيدروجيني التي ينبغي الانتباه إليها. فالدقيق الأقوى الأكثر غلوتينًا يمكن تخميره مدة أطول. ولإيجاد قيم الأس الهيدروجيني لخبزك أنصح بأخذ قياسات عدة أثناء صنع عجينك:

\begin{enumerate}
  \item قِس قيمة الأس الهيدروجيني لبادئتك الحامضة قبل استخدامها
  \item افحص الأس الهيدروجيني بعد مزج جميع المكونات
  \item افحص الأس الهيدروجيني قبل التقسيم والتشكيل المسبق
  \item افحص الأس الهيدروجيني قبل التشكيل
  \item افحص الأس الهيدروجيني لعجينك قبل التخمير النهائي وبعده
  \item افحص الأس الهيدروجيني لخبزك بعد الخَبز
\end{enumerate}

فإذا نجح الخبز الذي صنعته بهذه القيم يمكنك استخدامها مرجعًا للدفعة التالية. وإذا لم يخرج الخبز كما تحب، فإما قصّر التخمير أو مده قليلًا.

\begin{table}[htb]
    \centering
        % table-ph-values-dough-ar.tex — Arabic edition of upstream
% tables/table-ph-values-dough.tex (ch. 9, جدول 9.4)
%
% Original: The Sourdough Framework by Hendrik Kleinwächter,
% the-sourdough-framework.com, CC BY-SA 4.0. The Arabic edition inherits
% CC BY-SA 4.0.
%
% RTL mirror per the settled table mechanics (unit 9): the column order
% is reversed in the declared LTR tabular (upstream {@{}lr@{}} becomes
% {@{}rl@{}}) so the step label column renders rightmost and the pH
% value column leftmost, matching the accepted edition's render of
% p. 95. Every non-empty cell wrapped {\textdir TRT ...} per the unit 9
% cell-direction convention; the pH values ride \LR{} islands
% (siunitx output must stay LTR inside RTL paragraphs) and the pH
% header keeps the Latin symbol in an \LR{} island per the glossary
% (never translate the symbol). Cell text follows the accepted
% edition's render of p. 95 line for line; the caption lives in the
% chapter file.

\begin{tabular}{@{}rl@{}}
\toprule
\thead{{\textdir TRT المرحلة}} & \thead{{\textdir TRT قيمة \LR{pH}}} \\
\midrule
{\textdir TRT البادئة جاهزة}          & {\textdir TRT \LR{\num{4.20}}} \\
{\textdir TRT الخلط}                  & {\textdir TRT \LR{\num{6.00}}} \\
{\textdir TRT التقسيم والتشكيل المسبق} & {\textdir TRT \LR{\num{4.10}}} \\
{\textdir TRT التشكيل}                & {\textdir TRT \LR{\num{4.05}}} \\
{\textdir TRT قبل التخمير النهائي}     & {\textdir TRT \LR{\num{4.03}}} \\
{\textdir TRT بعد التخمير النهائي}     & {\textdir TRT \LR{\num{3.80}}} \\
{\textdir TRT بعد الخَبز}              & {\textdir TRT \LR{\num{3.90}}} \\
\bottomrule
\end{tabular}
%
        \caption[الأس الهيدروجيني للعجين أثناء تحضير الخبز]{قيم أس هيدروجيني نموذجية للنقاط الفاصلة المختلفة في عملية الخبز الحامض عندي.}
        \label{table:sample-ph-values}
\end{table}

وجمال هذه الطريقة موثوقيتها: فبمجرد إيجاد قيمك العاملة الجيدة يمكنك إعادة إنتاج مستوى التخمير نفسه مع كل عجين لاحق. وهذه مفيدة بادئ ذي بدء للمخبزات واسعة النطاق التي تريد تحقيق الاتساق في كل رغيف.

ورغم أن هذه الطريقة موثوقة جدًّا، فهناك قيود معينة يجب مراعاتها.

أولًا وقبل كل شيء، فإن قيم الأس الهيدروجيني التي تناسبني غالبًا لن تناسبك: فبحسب تركيبة بادئتك من بكتيريا حمض اللبنيك وحمض الخليك ستكون قيمك مختلفة. ويمكنك استخدام القيم المعروضة في الجدول~\ref{table:sample-ph-values} أرقامًا تقريبية فحسب، لكنك تحتاج على أي حال إلى إيجاد قيم تناسب إعداداتك.

والقيد الآخر هو السعر. فسيتعين عليك شراء مقياس أس هيدروجيني متطور، ومثاليًّا مقياس برأس سنّي.\footnote{ليس كل مقياس أس هيدروجيني مناسبًا لقياس العجين. راجع الدليل للتأكد من أنه معتمد لقياس الأس الهيدروجيني في الوسائط السائلة وشبه الصلبة. ولكي تحصل على قراءات دقيقة تأكد أيضًا من معايرة مقياسك معايرة سليمة.} وبهذا يمكنك غرز المقياس عميقًا في العجين مباشرةً. وفي الوقت نفسه ابحث عن ميزة الضبط التلقائي للحرارة؛ فبحسب درجة الحرارة تتغير قيمة الأس الهيدروجيني، وهناك جداول يمكنك استخدامها في حسابات الضبط، والمقاييس الأغلى تضم هذه الميزة مدمجة. ويفقد مقياس الأس الهيدروجيني دقته مع الوقت، ولهذا يلزمك معايرته باستمرار، وهي عملية شاقة تستغرق وقتًا. وأخيرًا عليك أن تغسل مقياس الأس الهيدروجيني بعناية قبل استخدامه في عجينك؛ فالسائل المحيط برأس المقياس غير صالح للأغذية ولا ينبغي أكله. وأنا أغسل المقياس مدة دقيقة على الأقل قبل استخدامه لقياس مرحلة تخمير عجيني.

والطريقة الأخيرة للحكم على حالة التخمير الأول هي قراءة علامات عجينك. فكلما صنعت خبزًا أكثر ستألف هذه العملية أكثر. راقب زيادة حجم العجين، وقد يكون هذا تحديًا أحيانًا حين يكون عجينك داخل وعاء. ويمكنك أن تساعد نفسك بوسم وعائك، بل يستخدم بعض الخبازين وعاء تخمير أول مستطيلًا شفافًا: يمكنك الوسم بقلم عند نقطة البداية الأولية ثم تراقب زيادة الحجم بوضوح. وكما في عينة الأليكوت المذكورة راقب زيادة حجم تناسب تركيبة بادئتك الحامضة.

\begin{figure}[htb]
  \centering
  \includegraphics[width=\textwidth]{bulk-finished-dough}
  \caption[عجين في نهاية التخمير الأول]{عجين في حالة جيدة لإنهاء التخمير الأول. لاحظ الفقاعات الدقيقة على سطح العجين؛ فهي علامة على أن العجين قد انتفخ بالغاز انتفاخًا كافيًا.}
\end{figure}

وراقب الفقاعات على سطح عجينك؛ فهي علامة جيدة على أن عجينك انتفخ بالغاز. وكلما دفعت التخمير الأول أبعد ظهرت فقاعات أكثر. وإذا أفرطت في هذه المرحلة أصبح العجين مسربًا واختفت الفقاعات من جديد.

ودوّن رائحة العجين؛ فينبغي أن توافق رائحة البادئة الناضجة قبل انهيارها بمدة قصيرة. وكما ذكرنا سابقًا، عجينك ليس إلا بادئة عملاقة. ويمكنك أيضًا أن تمضي وتذوّق عجينك: فسيذوق مثل المخللات. وبحسب الحموضة يمكنك الحكم على مدى تقدم العجين في عملية التخمير. وسيكون مذاق الخبز النهائي أقل حموضة، لأن كمية كبيرة من الحموضة تتبخر أثناء الخَبز.\footnote{المزيد عن هذا الموضوع لاحقًا. فبمجرد الخَبز مدة أطول أو أقصر يمكنك ضبط حموضة خبزك النهائي المخبوز: فكلما خبزت مدة أطول كان الرغيف النهائي أقل حموضة، وكلما خبزت أقصر بقيت حموضة أكثر داخل الخبز وكان الرغيف الناتج أكثر حموضة.}

وعند لمس العجين ينبغي أن يُستشعر ملتصقًا إلى يديك، وأن يكون أقل لزوجة مقارنة بالمراحل السابقة. فإذا كان العجين لزجًا أكثر من اللازم فقد دفعت التخمير أبعد من اللازم.

وإذا دفعت التخمير الأول أبعد من اللازم فلن تتمكن من خَبز رغيف قائم بذاته بهذا العجين بعد الآن. لكن لا تقلق: يمكنك نقل عجينك إلى قالب خبز، أو استخدام أجزاء من العجين بادئةً لعجينك التالي. وعند استخدام قالب الخبز احرص على تدهنه كما يجب. وقد تحتاج إلى سباتولا لنقل عجينك. ودع العجين يتخمر نهائيًّا مدة \LR{30} دقيقة على الأقل في القالب قبل خبزه؛ فهذا يضمن أن التجويفات الكبيرة التي سببها النقل تُسوّى. ويمكنك دفع التخمير النهائي إلى \LR{24} ساعة أو حتى \LR{72} ساعة، وسيتميز الخبز الناتج بمذاق رائع لاذع جدًّا.

\section{الطيّات والشدّات}

في هذا القسم ستتعلم كل ما تحتاج معرفته عن الشدّ والطيّ. ستتعلم متى تشدّ وتطوي، وكيف توظف هذه التقنية لصالحك.

\begin{figure}[htb]
  \centering
  \includegraphics[width=\textwidth]{dough-being-glued}
  \caption[لصق العجين]{عجين يُلصَق فيه سطحان لزجان أحدهما بالآخر بواسطة طيّة وشدّة. تمنح هذه العملية العجين قوةً ممتازة.}
\end{figure}

الشدّ والطيّ مجموعة تقنيات يستخدمها الخبازون أثناء مرحلة التخمير الأول. وتتضمن العملية شدّ العجين ثم طيّه على ذاته. بعض الوصفات تدعو إلى شدّة وطيّة واحدة، وبعضها إلى عدة طيّات وشدّات.

الهدف الأساسي من هذه التقنية هو إمداد عجينك بقوة عجين إضافية. كما هو موضح في الشكل~\ref{fig:dough-strength-sourdough} هناك طرق متعددة لصنع قوة العجين.\footnote{في الحقيقة رأيتُ وصفات \textenglish{no-knead} عديدة تدعو إلى عدم العجن الابتدائي، ثم تطبيق الطيّات والشدّات خلال التخمير الأول. والزمن اللازم لإنجاز كل الطيّات يوازي على الأرجح زمن العجن الابتدائي المطلوب.} إذا لم تعجن كثيرًا في البداية، يمكنك بلوغ المستوى نفسه من قوة العجين بتطبيق الطيّات والشدّات لاحقًا. وكلما زادت الطيّات والشدّات التي تجريها، زادت قوة العجين التي تضيفها إلى عجينك. وستكون النتيجة رغيفًا أجمل مظهرًا بانتفاخ فرني عمودي أكبر.

أحيانًا، إذا كان العجين قابلًا للمد على نحو كبير وترطيبه مرتفعًا جدًّا، كان الشدّ والطيّ ضرورة. وبدونهما سيفتقر العجين نفسه إلى قوة عجين كافية ولن ينتفخ في الفرن إطلاقًا.

ومن فوائد الطيّات والشدّات أيضًا خصائصها في التجنيس. فبطيّك العجين تعيد توزيع المناطق التي تتخمر أسرع من غيرها. فالحرارة في عجينك ليست واحدة في جميع المناطق، والتخمير ذاته يولّد حرارة، ولهذا ستتخمر بعض مناطق عجينك أسرع قليلًا من غيرها؛ ما يعني أن بعض المناطق تحتفظ بغاز وحمضية أكثر من غيرها. وتطبيق طيّة وشدّة يعيد توزيع الحرارة والغاز والحمضية. ويشير بعض الخبازين إلى هذه العملية باسم بناء اللُّبّ. والطيّات الحذرة تضمن ألا يكون لُبّ عجينك النهائي فوضويًا على نحو مفرط بتجاويف كبيرة. وإذا لاحظت تجاويف كبيرة على نحو مفرط في لُبّ عجينك النهائي، فقد تتمكن من إصلاح ذلك بتطبيق مزيد من الطيّات والشدّات.\footnote{في حالات كثيرة قد تحدث هذه التجاويف أيضًا حين لا يتخمر العجين ما يكفي. ويُسمّى اللُّبّ الناتج شائعًا لُبّ الأحمق \textenglish{Fool's Crumb}. راجع فصل تشخيص هياكل اللُّبّ لاحقًا في هذا الكتاب لمعرفة المزيد عنه.} راجع القسم~\ref{sec:debugging-crumb-structure} «تشخيص بنية اللُّبّ» لمزيد من المعلومات عن قراءة لُبّك.

\begin{figure}[htb]
  \centering
  \includegraphics[width=\textwidth]{stretch-and-fold-steps}
  \caption[خطوات الطّي والشد]{نظرة عامة على الخطوات اللازمة لأداء الطيّات والشدّات لعجائن القمح.}
  \label{figure:stretch-and-fold-steps}
\end{figure}

يكمن سبب شعبية هذه التقنية في كفاءتها. فبشدّك العجين إلى الخارج تزيد مساحة سطحه، ثم تطوي العجين فوق ذاته فتُلصَق مساحات كبيرة من العجين بعضها ببعض فعليًّا. تخيّل ورقة تضع عليها الغراء، ثم تطوي الورقة، فتلتصق مساحات كبيرة من الورقة ببعضها. وكرر العملية نفسها مع مزيد من الغراء حتى تُنشئ طبقات متعددة من الورق والغراء. وهذا عين ما يحدث لعجينك؛ فبقليلٍ جدًّا من الحركات تكون قد وضعت الغراء على عجينك.

لتطبيق طيّة وشدّة ابدأ بترطيب يديك بماء بارد، فالأيدي المبلولة تعمل العجائب في تقليص ميل العجين إلى الالتصاق بهما. ثم فكّ العجين بحرص من حواف وعاء التخمير الأول، وذلك بوضع يدك بحرص عند حافة العجين ودفعها إلى أسفل على جدران الوعاء. وحين تبلغ القاع اسحب العجين قليلًا إلى الداخل؛ فينبغي أن يبقى العجين مكانه لا يرجع إلى حافة الوعاء. وحاول أن تكون سريعًا ما أمكن في هذه الحركة؛ فكلما بطئت التصقت كمية أكبر من العجين بيديك. وكرر العملية نفسها حول العجين كله حتى يتخلص من حواف الوعاء. ثم بلّل يديك مرة أخرى وارفع جانبًا واحدًا من العجين إلى أعلى بيديك معًا موضوعتين في المنتصف، وأجرِ طيّة في منتصف العجين؛ فينبغي وضع السطح الأملس العلوي على قاع الوعاء. وبفعل ذلك تُلصَق الجانبان اللذان في الأسفل واللذان لزجان. أما السطح الأملس العلوي فلا ينبغي أن يكون لزجًا في يديك، بينما يميل السطح الخشن السفلي إلى الالتصاق بهما. وأدر الوعاء وكرر الأمر نفسه من الجهة الأخرى. وأدر الوعاء \LR{\ang{90}} ثم كرر العملية مرة أخرى. وأدر الوعاء \LR{\ang{180}} أخرى في الاتجاه نفسه وكرر الطيّة مرة أخيرة. وبفعل ذلك تكون قد طبّقت أربع طيّات في المجموع؛ فينبغي أن يبقى عجينك مكانه ويقاوم التدفق إلى الخارج.\footnote{راجع أيضًا~\cite{stretch+and+fold+technique} لفيديو يريك أفضل أداء لهذه التقنية.}

نظريًّا لا حدّ لعدد مرات الشدّ والطيّ؛ يمكنك تطبيق طيّة كل \LR{15} دقيقة. أما إذا كان عجينك يملك قوة عجين كافية أصلًا، فتطبيق طيّات إضافية مجرد إهدار للوقت.\footnote{يمكنك فعلها لمجرد أن تفهم أفضل شعور العجين في يديك عند مراحل تخمير مختلفة.} وإذا طبّقت عددًا كبيرًا من الطيّات المتتالية فسيتمزق طبقة الغلوتين الخارجية. وفي هذه الحالة عليك فقط انتظار \LR{\textenglish{5\textendash 10}} دقائق على الأقل حتى تلتئم روابط الغلوتين ثم تحاول مجددًا. وحين لا يلتئم الغلوتين بعد الآن فمن المرجح أنك دفعت التخمير مدةً أطول من اللازم، ومن المرجح أن معظم الغلوتين قد تفكك وأنك بلغت بالفعل مرحلة التدهور الموضحة في الشكل~\ref{fig:wheat-yeast-sourdough-degradation}.

\begin{figure}[htb]
  \centering
  \includegraphics[width=\textwidth]{dough-requiring-stretch-and-fold}
  \caption[عجين مفلطح]{عجين أثناء التخمير الأول قد مفلطح. ولتحسين قوته ينبغي تطبيق طيّة وشدّة.}
\end{figure}

أما الآن فكمية الطيّات والشدّات المعقولة التي ينبغي أن تجريها تتوقف بدرجة كبيرة على مقدار ما عجنته ابتداءً وعلى مدى قابلية عجينك للمد. والتوصية الجيدة أن تراقب عجينك في وعاء التخمير الأول، وحالما تراه يفلطح كثيرًا وينتشر نحو حواف الوعاء يمكنك المباشرة بتطبيق طيّة وشدّة. وبالنسبة إلى \LR{95\,\%} من العجائن التي أصنعها لا يتجاوز ذلك مرة واحدة تقريبًا. فأنا أحب صناعة عجائن تُترك لليل كامل، وعندها أطبق عادة طيّة وشدّة واحدة مباشرة بعد استيقاظي. ثم قد يستغرق التخمير الأول ساعتين إضافيتين قبل أن أشرع في التقسيم والتشكيل المسبق أو التشكيل مباشرة.

\section{اختياري: التقسيم والتشكيل المسبق}

التقسيم والتشكيل المسبق خطوة اختيارية تُنجز حالما ينتهي الخبز الحامض لديك من مرحلة التخمير الأول. والخطوة مطلوبة إذا كنت تصنع أرغفة متعددة في دفعة واحدة، واختيارية إذا كنت تصنع رغيفًا واحدًا.

\begin{flowchart}[htb]
  \centering
  % fig-dividing-preshaping-ar.tex — Arabic edition of upstream
% figures/fig-dividing-preshaping.tex (ch. 9, المخطط الانسيابي 9.5)
%
% Original: The Sourdough Framework by Hendrik Kleinwächter,
% the-sourdough-framework.com, CC BY-SA 4.0. The Arabic edition inherits
% CC BY-SA 4.0.
%
% R1-FIGURES route: re-derived from the upstream source
% figures/fig-dividing-preshaping.tex (node styles from
% figures/flowcharts_tikz.tex, node/edge geometry verbatim), with the
% label text recovered from the accepted edition's render of p. 101
% (book-ar-fixed.pdf, b78c2a28…). Line splits follow the accepted
% render (أكثر من/رغيف واحد؟). Arabic label text wrapped in
% \textarabic{}; the digit 1 rides an \LR{} island (Ruling 1
% uniform-Pagella figure face, ACCEPTED 2026-11-15).

\begin{tikzpicture}[node distance = 3cm, auto]
  \node [start] (init) {\textarabic{هل يلزم التقسيم؟}};
  \node [decision, right of=init, node distance=5cm] (more_than_one_loaf) {\textarabic{أكثر من}\\ \textarabic{رغيف واحد؟}};
  \node [success, right of=more_than_one_loaf, node distance=5cm] (yes) {\textarabic{نعم}};
  \node [success, below of=yes] (no) {\textarabic{لا}};
  \path [line] (init) -- (more_than_one_loaf);
  \path [line] (more_than_one_loaf) -- (yes);
  \path [line] (more_than_one_loaf.south) -- node{} ++(0, -1) |- (no);
\end{tikzpicture}

  \caption[التقسيم والتشكيل المسبق]{التقسيم مطلوب فقط حين تصنع أرغفة متعددة من دفعة عجين واحدة.}
  \label{fig:dividing-decision-tree}
\end{flowchart}

هدف تقسيم عجينك إلى قطع أصغر هو توزين عجينك وفقًا لذلك؛ فتحصل على قطع خبز متعددة تزن جميعها المقدار نفسه. ولهذا يُستخدم الميزان شائعًا لوزن قطع العجين. وإذا وازنت قطعة عجين أقل من اللازم يمكنك ببساطة قطع المزيد من كتلة عجينك لزيادة وزنها.

وعند قطع العجين حاول أن تكون حركاتك محكمة قدر الإمكان؛ فلا تريد إتلاف عجينك إتلافًا غير ضروري. والحركات السريعة بالسكين أو مكشطة العجين \textenglish{(dough scraper)} تساعد على منع العجين من الالتصاق بأدواتك كثيرًا.

\begin{figure}[htb]
  \centering
  \includegraphics[width=\textwidth]{divide-preshape}
  \caption[خطوات تقسيم عجينك وتشكيله مسبقًا]{خطوات تقسيم عجينك وتشكيله مسبقًا.}
\end{figure}

أحب أحيانًا رسم خطوط صغيرة بحافة مكشطة العجين على كتلة العجين الكبيرة قبل تقطيعها إلى قطع أصغر؛ فهذا يساعدني على تخطيط مواضع شقوقي أفضل. وحين أخطط لصنع \LR{8} أرغفة أحاول استخدام الخطوط لتقسيم العجين إلى \LR{8} أجزاء متساوية الحجم قبل التقطيع. وإذا لم يكن ذلك دقيقًا بما يكفي يمكنك استخدام الميزان المذكور آنفًا.

والآن بعد قطع عجينك، فالقطع الناتجة ليست في شكل متساوٍ، وهذا إشكالي للمرحلة التالية حين تشكّل عجينك؛ فلن تبدو الأرغفة الناتجة جميلة ومتساوية، وستنتهي على الأرجح بمناطق تمزق لحظة تشكيل عجينك، ولن تبدأ عملية التخمير النهائي كلها على أساس متين. ولهذا ينبغي تشكيل عجينك مسبقًا.

يُجرى التشكيل المسبق لعدة أسباب:

\begin{itemize}
  \item قسّمت عجينك وتحتاج إلى تشكيله مسبقًا
  \item عجينك يفتقر إلى قوة العجين، والتشكيل المسبق سيضيف مزيدًا من القوة
  \item تريد تسوية بنية اللُّبّ في الرغيف النهائي؛ فبالتشكيل المسبق سيبدو اللُّبّ الناتج أكثر انتساوية
\end{itemize}

وإذا كنت تصنع رغيفًا واحدًا من دفعة عجين واحدة فالخطوة غير مطلوبة، وعندها يمكنك المباشرة بالتشكيل مباشرة متجاوزًا هذه الخطوة.

وتقنية التشكيل المسبق هي نفسها العملية الموضحة في الشكل~\ref{fig:dough-ball-steps}. غير أن ما كان يمكنك فيه تمزيق سطح العجين سابقًا قد يصبح الآن كارثة. ولهذا أنصحك بالتدرب على هذه الخطوة ما دمت بحاجة إلى ذلك بعد العجن. فقد تكون شبكة الغلوتين عند هذه النقطة قابلة للمد ومتدهورة إلى حد لا يكاد يترك مجالًا لأي خطأ، ولن يعود العجين يتكامل مجددًا، والسبيل الوحيد لإنقاذ مثل هذا العجين هو قالب الخبز.

\begin{figure}[htb]
  \centering
  \includegraphics[width=\textwidth]{preshape-direction}
  \caption[اتجاه السحب]{اسحب العجين في اتجاه منطقة السطح الخشن. وبذلك تقلل الحركات اللازمة لإتمام الخطوة.}
  \label{fig:preshape-direction}
\end{figure}

شكّل عجينك مسبقًا بالقدر اللازم لتدوير مساحة السطح العلوي واستدارته. وحاول لمس العجين أقل ما يمكن لتقليل قدرته على الالتصاق بيديك، واسحب العجين في الاتجاه الذي ترى فيه منطقة سطح خشنة. وإذا ضاقت بك المساحة للسحب لأن العجين قد يسقط من حافة رخامتك فببساطة ارفعه بحركة سريعة وضعه في موضع أفضل للتشكيل المسبق. راجع الشكل~\ref{fig:preshape-direction} لتمثيل بصري يوضح اتجاه التشكيل المسبق.

وحاول أن تضع لنفسك حدًّا من الحركات لإتمام تشكيل العجين مسبقًا، وستكون حينها أكثر وعيًا بكل حركة تجريها. في البداية يمكنك تجربة \LR{5} حركات ثم تخفيضها تدرّجيًّا إلى \LR{3}. والسبب الوحيد لتجاوز هذه الأرقام قد يكون رغبتك المقصودة في تسوية بنية اللُّبّ في أرغفتك النهائية مزيدًا.

\begin{figure}[htb]
  \centering
  \includegraphics[width=\textwidth]{preshaped-dough}
  \caption[عجائن باجيت ترتاح بعد التشكيل المسبق]{عجائن باجيت ترتاح بعد التشكيل المسبق.}
  \label{fig:dough-after-preshaping}
\end{figure}

حين تنتهي من التشكيل المسبق اترك كرات العجين ترتاح على رخامتك \LR{\textenglish{10\textendash 15}} دقيقة على الأقل. ولا تغطِّ الكرات المشكَّلة مسبقًا؛ فتجفيف السطح يجعل خطوة التشكيل التالية أسهل، ولن يلتصق السطح المجفف بيديك كثيرًا. وبما أنك شددت غلوتين العجين فستحتاج إلى إتاحته للاسترخاء. وبدون فترة راحة لن تتمكن من تشكيل عجينك في مثل بنية الباجيت مثلًا؛ فسيقاوم العجين كل حركة عائدًا دائمًا إلى الشكل السابق. ولعلك لاحظت هذا من قبل عند صناعة عجين البيتزا؛ فإذا لم تنتظر ما يكفي بعد تكوير البيتزات استحال مدّ البيتزا، وبانتظار بضع دقائق إضافية يصبح المدّ أسهل بكثير، ولن يقاوم العجين التحول إلى الشكل النهائي الذي تريده.

وزمن الراحة على الرخامة المذكور آنفًا والبالغ \LR{\textenglish{10\textendash 15}} دقيقة يتوقف على مدى قوة تشكيلك المسبق لعجينك؛ فكلما شكّلته مسبقًا أكثر احتجت إلى انتظار أطول. وإذا قاوم عجينك كثيرًا أثناء التشكيل فمُد هذه الفترة حتى \LR{30} دقيقة. وإذا انتظرت طويلًا أكثر من اللازم فقد يصبح سطح عجينك جافًا أكثر من اللازم، فيتمزق العجين أثناء التشكيل. وكالمعتاد خذ هذه الأزمنة بحبة ملح وجرب في بيئتك.

\section{التشكيل}

\begin{flowchart}[htb]
\centering
% fig-shaping-process-ar.tex — Arabic edition of upstream
% figures/fig-shaping-process.tex (ch. 9, المخطط الانسيابي 9.6)
%
% Original: The Sourdough Framework by Hendrik Kleinwächter,
% the-sourdough-framework.com, CC BY-SA 4.0. The Arabic edition inherits
% CC BY-SA 4.0.
%
% R1-FIGURES route: re-derived from the upstream source
% figures/fig-shaping-process.tex (node styles from
% figures/flowcharts_tikz.tex, node/edge geometry verbatim), with the
% label text recovered from the accepted edition's render of p. 104
% (book-ar-fixed.pdf, b78c2a28…). Line splits follow the accepted
% render (هل العجين لزجًا/أكثر من اللازم/أو يتمزق؟، عجينك على الأرجح/
% مفرط التخمير، انقل إلى قالب/الخبز، ثم تخمير/نهائي قصير،/واخبز مباشرة،
% تابع بتقنية/التشكيل، رشّ العجين/المشكَّل بالدقيق، ضعه مقلوبًا/
% في سلة التخمير). Arabic label text wrapped in \textarabic{} (Ruling 1
% uniform-Pagella figure face, ACCEPTED 2026-11-15). Decision edge
% labels نعم/لا follow the accepted render. The wide decision node
% keeps the standard decision style; if the build shows the unit 12
% line-break problem, the local decisionnat style is the settled
% fallback (unit 12 convention).

\begin{tikzpicture}[node distance = 3cm, auto]
  \node [start] (init) {\textarabic{ابدأ التشكيل}};
  \node [decision, right of=init, node distance=4cm] (overfermented_decision) {\textarabic{هل العجين لزجًا}\\ \textarabic{أكثر من اللازم}\\ \textarabic{أو يتمزق؟}};
  \node [block, right of=overfermented_decision, node distance=4cm] (overfermented) {\textarabic{عجينك على الأرجح}\\ \textarabic{مفرط التخمير}};
  \node [fail, right of=overfermented, node distance=4cm] (loafpan) {\textarabic{انقل إلى قالب}\\ \textarabic{الخبز، ثم تخمير}\\ \textarabic{نهائي قصير،}\\ \textarabic{واخبز مباشرة}};
  \node [block, below of=overfermented_decision, node distance=4cm] (shaping_technique) {\textarabic{تابع بتقنية}\\ \textarabic{التشكيل}};
  \node [block, right of=shaping_technique] (flour) {\textarabic{رشّ العجين}\\ \textarabic{المشكَّل بالدقيق}};
  \node [block, right of=flour] (banneton) {\textarabic{ضعه مقلوبًا}\\ \textarabic{في سلة التخمير}};
  \node [success, right of=banneton] (proof) {\textarabic{ابدأ التخمير النهائي}};
  \path [line] (init) -- (overfermented_decision);
  \path [line] (overfermented_decision) -- node{\textarabic{نعم}} (overfermented);
  \path [line] (overfermented_decision) -- node{\textarabic{لا}} (shaping_technique);
  \path [line] (shaping_technique) -- (flour);
  \path [line] (flour) -- (banneton);
  \path [line] (banneton) -- (proof);
  \path [line] (overfermented) -- (loafpan);
\end{tikzpicture}

\caption[عملية التشكيل]{تمثيل تخطيطي لعملية التشكيل متضمنًا فحوصات عجين مفرط التخمير.}
\label{fig:shaping-decision-tree}
\end{flowchart}

التشكيل يمنح عجينك الشكل النهائي قبل الخَبز. وبعد إتمام التشكيل ينتقل عجينك إلى مرحلة التخمير النهائي ثم يُشرَّط وفي النهاية يُخبَز.

تقنيات التشكيل لا تُحصى، والتقنية التي تختارها تتوقف على نوع الخبز الذي تريد صناعته. بعض التقنيات ألطف على العجين وتضمن ألا يفقد العجين غازه، وأخرى أسرع لكنها تُخرِج الغاز منه قليلًا أكثر. وكلما شُدَّت بإحكام أكبر بدت بنية اللُّبّ في عجينك النهائي أكثر انتساوية، وفي الوقت نفسه سيحسّن التشكيل الأكثر إحكامًا قوة عجينك، والمزيد من القوة سينتج في النهاية انتفاخًا فرنيًا عموديًا أكبر.

يفترض ما يأتي من تعليمات أنك تريد صنع خبز بنمط \textenglish{batard} ذاتي الشكل الممتد الممدود. وتعلم هذه التقنية سيمنحك أساسًا معرفيًا متينًا يمكن توسيعه بسهولة لصنع لفائف الخبز أو الباجيتات.

وإتقان تقنية التشكيل الصعبة سيستغرق منك على الأرجح محاولات متعددة، غير أن أمامك محاولة واحدة فقط لكل عجين؛ فإذا أخطأت فلن يخرج الخبز النهائي على الأرجح بأفضل ما يمكن. وإذا سبّبت لك هذه التقنية صداعًا فأنصحك بصنع دفعة عجين أكبر وتقسيمها وتشكيلها مسبقًا إلى حصص أصغر، وبدل صنع باتارد كبير تدرّب على صنع لفائف باتارد مصغّرة.

\subsection[نثر الدقيق على السطح]{افرد الدقيق على سطح العجين.}

إذا كنت تصنع رغيفًا واحدًا فقط من عجينك فافرد الدقيق بسخاء على الطبقة العلوية من عجينك، وافرك الدقيق على عجينك بيديك، ثم اقلب وعاءك وانتظر قليلًا ليتحرر العجين من الوعاء، ثم المباشرة بالخطوة \LR{3}.

وإذا قسّمت عجينك وشكّلته مسبقًا فافرد الدقيق بسخاء على الطبقة العليا للعجين كذلك، وبأيدي رقيقة انشر الدقيق بانتساوية عبر سطح العجين. راجع الشكل~\ref{fig:shaping-flour-surface} لتمثيل بصري لما يجب أن يبدو عليه عجينك بعد طلاء السطح.

\begin{figure}[htb]
  \centering
  \includegraphics[width=0.7\textwidth]{step-1-flour-applied}
  \caption[الخطوة 1 من عملية التشكيل]{عجين وُضِع الدقيق على سطحه. هذه هي الخطوة الأولى من عملية التشكيل.}
  \label{fig:shaping-flour-surface}
\end{figure}

\subsection[قلب العجين]{اقلب العجين}

بأيدي رقيقة انزع العجين بحرص من السطح. وإذا كان لديك مكشطة عجين فأدخلها تحت العجين بحرص بحركات سريعة. واقلب العجين مضمنًا أن مناطق الدقيق المفروك ملامسة ليديك. أما المنطقة السفلية غير المفروكة التي كانت ملتصقة بالرخامة فمنطقة محرَّمة اللمس؛ فحاول تجنب لمسها لأنها خشنة وبالتالي ستلتصق بيديك.

ثم ضع العجين بلطف بحيث يواجه الرخامة الجانب الذي كان يواجه الأعلى سابقًا؛ فتكون منطقة الدقيق المفروك على السطح الآن بينما يواجهك الجانب اللزج.

\begin{figure}[htb]
  \centering
  \includegraphics[width=\textwidth]{step-2-flipped-over}
  \caption[الخطوة 2 من عملية التشكيل]{عجين مقلوب. لاحظ كيف يواجهك الجانب اللزج بينما يواجه الدقيق المفروك سطح الرخامة. يُستخدم الجانب اللزج غراءً لإمساك العجين بعضه ببعض.}
  \label{fig:shaping-rectangular-dough}
\end{figure}

\subsection[إعطاء العجين شكلًا مستطيلًا]{اجعل العجين مستطيلًا}

ينبغي أن يكون الجانب اللزج من عجينك مواجهًا لك الآن. ولاحظ كيف أن العجين مستدير حاليًا وليس مستطيلًا، والشكل المستدير لن يكون مثاليًا عند تشكيل الباتارد الممتد الممدود.

ولهذا اشدد العجين قليلًا حتى يأخذ شكلًا أكثر استطالة. وأثناء الشدّ احرص على لمس الجانب اللزج أقل ما يمكن. ضع يديك على الجانب السفلي المفروك وعلى حافة الجانب اللزج، وبأيدي رقيقة اشدد العجين حتى يبدو الشكل أمامك مستطيلًا. راجع الشكل~\ref{fig:shaping-rectangular-dough} وقارن عجينك بالعجين الموضح.

\begin{figure}[htb]
  \centering
  \includegraphics[width=\textwidth]{step-3-rectangular}
  \caption[الخطوة 3 من عملية التشكيل]{عجين مقلوب. لاحظ كيف يواجهك الجانب اللزج بينما يواجه الدقيق المفروك سطح الرخامة.}
  \label{fig:shaping-folding}
\end{figure}

\subsection[الطّيّ]{اطوِ العجين بعضه على بعض}

والآن بعد أن أنشأت الشكل المستطيل فعجينك جاهز ليُطوى بعضه على بعض. وهذا لا ينجح إلا لأن الجانب المواجه لك لزج، وبفضل لزوجة العجين يمكننا لصقه بعضه ببعض بفعالية فينشأ رابط قوي جدًا.

يمكنك التدرب على هذه الخطوة بورقة مستطيلة، وحين تتقن الطّيّ على الورق يمكنك بسهولة تطبيقها على عجينك الحقيقي.

احرص على أن يكون الباتارد موضوعًا أمامك. خذ الجانب المواجه لك واطوِه نحو منتصف العجين، وادّسه بحرص لأسفل حتى يلتصق بالجانب اللزج.

\begin{figure}[htb]
  \centering
  \includegraphics[width=\textwidth]{step-4-folding}
  \caption[الخطوة 4 من عملية التشكيل]{عملية طيّ الباتارد. لاحظ كيف يُلصَق المستطيل أولًا بعضه ببعض ثم يُلف إلى الداخل لصنع لفافة عجين، وتُغلق الحواف أخيرًا لصنع عجين أكثر تجانسًا.}
\end{figure}

خذ الجانب الآخر واطوِه فوق الجانب الذي طويته للتو. واشدد العجين نحوك أقصى ما يمكن، وادّسه عند الحافة لتنشئ طبقة الغراء الأولى.

أدر العجين حتى يستقيم طوله أمامك، وأدره إلى الداخل بحيث يواجهك جانب الدرزة الآن.

ابدأ بلّف العجين إلى الداخل بدءًا من أعلى العجين، وواصل لّفه إلى الداخل حتى تنشئ لفافة عجين.

راجع الشكل~\ref{fig:shaping-folding} لتمثيل بصري كامل للعملية.

وإذا لم يحتفظ عجينك بشكله فمن المرجح أنك دفعت التخمير بعيدًا أكثر من اللازم؛ فقد تفكك معظم الغلوتين ولن يقدر العجين على الاحتفاظ بشكله. وعندها فالخيار الأفضل استخدام قالب الخبز لخبز خبزك؛ فسيكون طعم الخبز النهائي مذهلًا لكنه لن يوفر القوام نفسه الذي يوفره الخبز القائم بذاته. راجع القسم~\ref{sec:debugging-crumb-structure} لمزيد من التفاصيل عن كيفية قراءة بنية اللُّبّ في عجينك قراءة صحيحة.

\subsection{غلق الحواف}

عجينك أنهى التشكيل الآن. وغلق الحواف خطوة اختيارية، وأحب القيام بها لأنني أرى أن الخبز النهائي المخبوز سيبدو أجمل قليلًا بلا حواف خشنة.

اسحب حواف عجينك الحلزونية المظهر ببلطف بإصبعين، وأدر العجين ثم كرر العملية نفسها من الجهة الأخرى كذلك.

\subsection[الاستعداد للتخمير النهائي]{استعد للتخمير النهائي}

\begin{figure}[htb]
  \centering
  \includegraphics[width=\textwidth]{step-6-prepare-proofing}
  \caption[الخطوة 5 من عملية التشكيل]{العجين المشكَّل جاهز للتخمير النهائي في سلة التخمير. لاحظ كيف يواجهك جانب الدرزة الآن، بينما الجانب العلوي المفروك السابق يواجه الأسفل.}
  \label{fig:shaping-prepare-proofing}
\end{figure}

ينبغي أن يكون أمامك الآن عجين جميل الشكل. ومرحلة التخمير النهائي على وشك البدء. ولتبسيط التشريط لاحقًا ولضمان ألا يلتصق عجينك بسلة التخمير، افرك سطح العجين بجوفة دقيق أخرى؛ فسيجفف هذا السطح ويقلل ميل العجين إلى الالتصاق بكل شيء.

وللطلاء أنصح باستخدام الدقيق نفسه الذي صنعته منه عجينك. والدقيق الأرز لا يوصى به إلا إذا أردت تطبيق أنماط تشريط فنية لاحقًا. ومن الأفضل استخدام دقيق أكثر من اللازم؛ فالدقيق الزائد يمكن تنظيفه بالفرشاة لاحقًا.

حالما يُطلَّى عجينك فهو جاهز ليوضع في سلة التخمير. وإذا لم تملك سلة تخمير فيمكنك استخدام وعاء بمنشفة مطبخ بداخله.

والمطحون المفروك المواجه للأعلى حاليًا سيواجه الأسفل في سلة التخمير. وبذلك يمكن قلب سلة التخمير قبل الخبز فيتحرر العجين.\footnote{ينطبق الأمر نفسه عند صناعة عجائن أخرى كعجائن الباجيت؛ فالسطح المفروك سيواجه الأسفل دائمًا، ويُقلب العجين مرة واحدة للخبز.}

ثم ارفع العجين بيديك معًا من الرخامة، وأدره بلطف مرة واحدة ثم ضعه في سلة التخمير للتخمير النهائي.\footnote{ينبغي أن يواجهك جانب الدرزة الآن. ويحب بعض الخبازين غلق الدرزة قليلًا أكثر، ولم ألاحظ أن ذلك يحسّن قوة العجين؛ فما أستطيع تمييزه أن هذا يحسّن المظهر البصري للجانب السفلي من الرغيف النهائي فقط.}

وإذا فعلت كل شيء صحيحًا فينبغي أن يشبه عجينك العجين الموضح في الشكل~\ref{fig:shaping-prepare-proofing} إلى حد بعيد. وخطوة أخيرة في التشكيل: ضع منشفة مطبخ على سلة التخمير أو الوعاء وابدأ التخمير النهائي.

\section{التخمير النهائي}

في صناعة الخبز يشير التخمير النهائي إلى النفخ الأخير للعجين قبل الخبز، بعد تشكيله رغيفًا. والتفاعلات والعمليات الكيميائية التي تجري أثناء التخمير الأول والتخمير النهائي واحدة.

فبتشكيلك عجينك فقد بعض الهواء المتولد سابقًا على مدار التخمير الأول، وهدف التخمير النهائي نفخ العجين من جديد. فعجين بلا تخمير نهائي لن يوفر القوام نفسه الذي يوفره عجين مخمَّر نهائيًا على نحو صحيح؛ فالعجين المخمَّر النهائي يتميز بلُبّ منفوش وطرلل جدًا.

وهناك تقنيتان للتخمير النهائي: إحداهما تخمير العجين في حرارة الغرفة، والأخرى تخميره في البراد. والتخمير النهائي في البراد معروف شائعًا باسم \textenglish{retarding} (الإبطاء).

يدّعي بعض الخبازين أن التخمير النهائي البارد يحسّن النكهة النهائية للخبز. وفي جميع الأرغفة التي أبطأتها لم أستطع تمييز فرق في النكهة بين العجائن المخمَّرة باردًا. فالكائنات الدقيقة تعمل بمعدل أبطأ عند الحرارات الأبرد، لكنني أشك أنهم يغيرون عملياتهم الكيميائية الحيوية. وما زال البحث مطلوبًا في موضوع الإبطاء وتطور النكهة.

\begin{flowchart}[htb]
\centering
% fig-proofing-process-ar.tex — Arabic edition of upstream
% figures/fig-proofing-process.tex (ch. 9, المخطط الانسيابي 9.7)
%
% Original: The Sourdough Framework by Hendrik Kleinwächter,
% the-sourdough-framework.com, CC BY-SA 4.0. The Arabic edition inherits
% CC BY-SA 4.0.
%
% R1-FIGURES route: re-derived from the upstream source
% figures/fig-proofing-process.tex (node styles from
% figures/flowcharts_tikz.tex, node/edge geometry verbatim), with the
% label text recovered from the accepted edition's render of p. 111
% (book-ar-fixed.pdf, b78c2a28…). ACCEPTED-EDITION DIVERGENCE,
% documented: the upstream node "Poke the dough" reads «انتظر 15 دقيقة»
% in the accepted edition — the accepted render's text layer carries
% three انتظر/15 دقيقة nodes (wait_poke far left, the poke position
% below init, wait_retard right) and no poke-test node, in both the raw
% and the layout-mode extraction. The accepted edition is the fidelity
% authority, so the node text follows the accepted render; flagged for
% render verification in the unit battery. Line splits follow the
% accepted render (التخمير النهائي/في حرارة/الغرفة؟، هل تخبز خلال/
% 10 ساعات/من الآن؟، انتظر/15 دقيقة، هل تبقى الحفرة/ظاهرة بعد/
% دقيقة واحدة؟، شقّ السطح/واخبز، التخمير في الثلاجة/والخبز عند
% 4°C (40 °F)). Arabic label text wrapped in \textarabic{}; digits,
% the pH-free temperature pair and °C/°F ride \LR{} islands (Ruling 1
% uniform-Pagella figure face, ACCEPTED 2026-11-15). Decision edge
% labels نعم/لا follow the accepted render.

\begin{tikzpicture}[node distance = 3cm, auto]
  \node [decision_start] (init) {\textarabic{التخمير النهائي}\\ \textarabic{في حرارة}\\ \textarabic{الغرفة؟}};
  \node [decision, right of=init, node distance=8cm] (retard_bake_decision) {\textarabic{هل تخبز خلال}\\ \LR{10} \textarabic{ساعات}\\ \textarabic{من الآن؟}};
  \node [block, below of=init, node distance=4cm] (poke) {\textarabic{انتظر}\\ \LR{15} \textarabic{دقيقة}};
  \node [block, left of=poke] (wait_poke) {\textarabic{انتظر}\\ \LR{15} \textarabic{دقيقة}};
  \node [decision, below of=poke] (dent_visible_decision) {\textarabic{هل تبقى الحفرة}\\ \textarabic{ظاهرة بعد}\\ \textarabic{دقيقة واحدة؟}};
  \node [success, right of=dent_visible_decision, node distance=4cm] (bake) {\textarabic{شقّ السطح}\\ \textarabic{واخبز}};
  \node [block] at (retard_bake_decision |- poke) (wait_retard) {\textarabic{انتظر}\\ \LR{15} \textarabic{دقيقة}};
  \node [block] at (wait_retard |- bake) (retard) {\textarabic{التخمير في الثلاجة}\\ \textarabic{والخبز عند \LR{4 °C (40 °F)}}};
  \path [line] (init) -- node{\textarabic{نعم}} (poke);
  \path [line] (init) -- node{\textarabic{لا}} (retard_bake_decision);
  \path [line] (poke) -- (dent_visible_decision);
  \path [line] (dent_visible_decision) -- node{\textarabic{نعم}} (bake);
  \path [line] (dent_visible_decision.west) -- node{\textarabic{لا}} (dent_visible_decision.west -| wait_poke.south) -- node{} (wait_poke.south);
  \path [line] (wait_poke) -- (poke);
  \path [line] (retard_bake_decision) -- node{\textarabic{نعم}} (wait_retard);
  \path [line] (retard_bake_decision.east) -- node{\textarabic{لا}} ++(1, 0) |- node{} (retard.east);
  \path [line] (wait_retard) -- (retard);
  \path [line] (retard) -- (bake);
\end{tikzpicture}

\caption[عملية التخمير النهائي]{نظرة تخطيطية عامة على الخطوات المختلفة لعملية التخمير النهائي للخبز الحامض. وتقنية التخمير النهائي التي تختارها تتوقف على جدولك ومواصلك.}
\label{fig:proofing-process}
\end{flowchart}

وبالنسبة إليّ فإن الغرض الوحيد من التخمير النهائي البارد هو قدرته على تمكينك من إدارة توقيت العملية كلها أفضل. فافترض أنك أنهيت تشكيل عجينك في الساعة \LR{10} مساءً، فلن ترغب على الأرجح في انتظار ساعتين إضافيتين لتخمير العجين ثم ساعة أخرى لخبزه. وفي هذه الحالة يمكنك نقل عجينك مباشرة إلى البراد بعد التشكيل، فيتخمر عجينك طوال الليل في البراد، ثم يمكن خبزه في أي وقت من اليوم التالي (ثمة استثناءات قليلة، والمزيد عن ذلك لاحقًا). وهذا مفيد بدرجة خاصة للمخابز واسعة النطاق التي تستخدم التخمير النهائي في البراد على نحو واسع؛ فبعض العجائن تُخمَّر قبل يوم وتوضع في البراد، وفي الصباح الباكر يمكن خبزها مباشرة من البراد، وخلال ساعتين ستكون جاهزة لبيع أول خبز لزبائن الصباح. وإذا احتاج المزيد من الخبز خلال اليوم فببساطة يأخذون بعض العجين المخمَّر من البراد ويخبزونه. والإطار الزمني الذي يمكنك فيه خبز العجين المُبطأ واسع، فيمكن أن يكون بعد \LR{6} ساعات فحسب وحتى بعد \LR{24} ساعة.

وافترض أنك صنعت عجين ليلة كاملة وعجينك جاهز في الصباح، فقد يختلف الموقف؛ فقد ترغب في خبز العجين مباشرة للفطور أو وقت الغداء، وعندها لن ترغب في تخمير العجين \LR{6} ساعات إضافية في البراد، والتخمير النهائي في حرارة الغرفة هو تقنيتك المفضلة.

وباختصار اختر التقنية التي تناسبك وفقًا لجدولك ومواصلك.

\subsection{التخمير النهائي في حرارة الغرفة}

أسهل الطرق وأكثرها موثوقية لتخمير عجينك نهائيًا هو تخميره في حرارة الغرفة. وهي طريقتي المفضلة إذا سمح لي جدولي. وتعمل هذه الطريقة بشكل رائع إذا صنعت عجين ليلة كاملة ثم خمّرته نهائيًا في صباح اليوم التالي.

\begin{figure}[htb]
  \centering
  \includegraphics[width=\textwidth]{step-13-finger-poke-test}
  \caption[اختبار النقر بالإصبع]{اختبار النقر بالإصبع \textenglish{(finger poke test)} طريقة موثوقة جدًا للتحقق من تخمير عجينك النهائي على نحو صحيح. فإذا بقيت الحفرة المحدثة ظاهرة بعد دقيقة كان عجينك صالحًا للخبز.}
  \label{fig:shaping-finger-poke}
\end{figure}

والزمن الذي يستغرقه تخمير عجينك نهائيًا قد يتراوح بين \LR{30} دقيقة و\LR{3} ساعات. وبدلًا من الاعتماد على التوقيت يستخدم معظم الخبازين اختبار النقر بالأصبع.

افرك إبهامك بالدقيق واضغط بلطف على عمق نحو \LR{\textenglish{0.5 cm\textendash 1 cm}} في العجين. وجرب هذا مباشرة بعد التشكيل وستلاحظ أن الحفرة المحدثة تتعافى سريعًا، فيختفي من جديد بعد دقيقة.

وبمضيك في التخمير النهائي سيمتلئ عجينك بمزيد من الغاز، وفي الوقت نفسه سيصبح العجين أكثر قابلية للمد. وحالما يبدأ في بلوغ الكمية المناسبة من الانفوش وقابلية المد ستتعافى الحفرة ببطء أكبر. وحالما يصبح العجين جاهزًا للتشريط والخبز ينبغي أن تبقى الحفرة ظاهرة بعد دقيقة من الانتظار.

أنصح بإجراء اختبار النقر بالإصبع مرة كل \LR{15} دقيقة على مدار مرحلة التخمير النهائي. وواقعيًّا، بحسب خبرتي، يستغرق التخمير النهائي ساعة على الأقل وقد يستغرق أحيانًا حتى \LR{3} ساعات. وحتى عند الحرارات الأدفأ لم يسبق أن كان التخمير النهائي أسرع من ساعة عندي. وكالمعتاد خذ أزمتي بحبة ملح وجرب بنفسك.

وحالما أرى أن العجين يقترب من التخمير النهائي المثالي أمضي وأسخّن فرني؛ فبذلك لا أفرط في تخمير العجين. وستلاحظ عجينًا مفرط التخمير حين يصبح العجين فجأة لزجًا جدًا، وفي الوقت نفسه من المرجح أن ينهار أثناء الخبز ولن يرتد. وعمومًا من الأفضل إنهاء التخمير النهائي مبكرًا أكثر مما ينبغي لا متأخرًا أكثر مما ينبغي.

\subsection{التخمير النهائي البارد \textenglish{(}الإبطاء\textenglish{)}}

الخيار الثاني للتخمير النهائي هو وضع عجينك داخل البراد لتخميره. وهذا الخيار رائع إذا لم تكن تريد خبز العجين خلال الساعات الثلاث التالية.

سيختمر العجين ابتداءً بالمعدل نفسه الذي يختمر به عجين حرارة الغرفة. وكلما برد العجين بطأ معدل التخمير، وفي النهاية عند درجة أقل من \LR{\qty{4}{\degreeCelsius}} (\LR{\qty{40}{\degF}}) يتوقف التخمير.\footnote{الحرارة الفعلية تتوقف على البكتيريا والخمائر التي ربيتها في بادئتك الحامضة.} ويمكن للعجين أن يرتاح في البراد حتى \LR{24} ساعة، وفي بعض التجارب بقي العجين جيدًا حتى بعد \LR{48} ساعة. والمثير للاهتمام أن هناك حدًا لتخمير البراد، ولا أستطيع تفسير ذلك إلا بنشاط تخميري متواصل عند الحرارات المنخفضة.

والجزء الصعب الحكم على متى انتهى عجينك من التخمير النهائي في برادك. فاختبار النقر بالإصبع المذكور سابقًا لا يعمل على العجين البارد؛ فالحرارات المنخفضة تغيّر مرونة العجين، وحفرة اختبار النقر لن تتعافى أبدًا.

ولهذا فإيجاد أفضل زمن للتخمير النهائي في البراد يُنجَز أفضل بما هو نهج تكراري. ابدأ بـ\LR{8} ساعات على عجينك الأول، و\LR{10} ساعات على الثاني، و\LR{12} ساعة على الثالث، وهكذا حتى \LR{24} ساعة. وبما أن الحرارة في برادك ثابتة عادة فأنت تملك بيئة يمكنك فيها الاعتماد على التوقيتات. فجد زمن التخمير النهائي المثالي الذي يناسبك.

واعتبار إضافي هو الحرارة الجوهرية للعجين قبل وضعه في البراد؛ فكلما كان عجينك أدفأ ابتداءً استغرق الأمر أطول ليبرد العجين، وهذا متغير إضافي ينبغي أخذه في الحسبان عند اختيار زمن الإبطاء. وفي الصيف حين يكون مطبخي حارًا أختار زمن تخمير نهائي في البراد أقصر مقارنة بالشتاء حين يكون العجين أبرد.

وطريقة موثوقة لضمان تخمير نهائي منتظم هي اختيار استخدام مقياس الأس الهيدروجيني. فبفحص كمية الحمضية المتراكمة يمكنك ضمان أن كل عجائنك تملك الكمية المناسبة من الحمضية. واختر نهجًا تكراريًا وافحص الأس الهيدروجيني لأزمنة تخمير نهائي متعددة، وابحث عن قيمة الأس الهيدروجيني التي تصنع أفضل خبز لك. وحالما تحدد قيمة الأس الهيدروجيني المثالية لك يمكنك الرجوع إلى ذلك الرقم في جميع العجائن التالية. راجع الجدول~\ref{table:sample-ph-values} لبعض قيم الأس الهيدروجيني النموذجية التي يمكن اتباعها.

\section{التشريط}

حالما ينتهي عجينك من التخمير النهائي حان وقت تسخين فرنك إلى نحو \LR{\qty{230}{\degreeCelsius}} (\LR{\qty{446}{\degF}}). والخطوة التالية هي حينها المباشرة بتشريط عجينك.

يُجرى التشريط لسببين: ثمة تشريط وظيفي وتشريط زخرفي. ويشير التشريط الوظيفي إلى إحداث شق صغير في العجين ينتفخ من خلاله أثناء الخبز. فإذا لم يُشرَّط العجين فمن المرجح أنه سيتشقق عند أضعف المواضع حيث أغلقت العجين بعد التشكيل. ويمكن استخدام التشريط الزخرفي لتطبيق أنماط فنية على عجينك وجعله أكثر جاذبية. وحين تريد تطبيق تشريط فني فمن الأفضل أن تفرك عجينك بدقيق أرز إضافي قبل التشريط؛ فالدقيق الأرز الأبيض يعزز تباين شقوق التشريط بدرجة كبيرة، وبالتالي يجعل النمط النهائي أكثر جاذبية بصريًا.

\begin{figure}[htb]
  \centering
  \includegraphics[width=\textwidth]{the-ear}
  \caption[أذن الخبز]{الأذن خاصية يمكن بلوغها في الخبز الحامض القمحي عند تخمير عجينك وتشريطه بالتقنية المثالية. وهي توفر نكهة إضافية وقوامًا رائعًا عند أكل الخبز.}
  \label{fig:the-ear}
\end{figure}

عند استخدام سلة تخمير يُقلب العجين ويوضع على رف فرن أو صينية أو حجر أو صلب أو وعاء خبز هولندي \textenglish{(dutch oven)}. ومحاسن وعيوب خيارات الخبز المختلفة يغطيها الفصل التالي. وسيواجهك الآن الجانب العلوي للعجين الذي كان سابقًا في قاع سلة التخمير.

\begin{figure}[htb]
  \centering
  \includegraphics[width=\textwidth]{artistic-scoring}
  \caption[تشريط فني]{رغيف لنانسي آن بنمط تشريط فني. وقد تحقق التباين العالي بفرك سطح العجين بدقيق الأرز قبل الخبز. وحسابها في إنستغرام \textenglish{\texttt{simply.beautiful.sourdough}} متخصص في عرض أنماط التشريط الفنية الجميلة.}
  \label{fig:artistic-scoring}
\end{figure}

ويُجرى شق التشريط بزاوية \LR{\ang{45}} بالنسبة إلى سطح العجين قليلًا بعيدًا عن مركز العجين. وبالشق بزاوية \LR{\ang{45}} سينتفخ الجانب الملصَق فوقه في الفرن أكثر من الجانب الآخر، وبذلك ستحقق ما يسمى الأذن في الخبز النهائي. والأذن حافة رقيقة مقرمشة توفر قوامًا شيقًا عند الأكل، وهي حافة رقيقة تكون عادة أغمق قليلًا بعد الخبز وبالتالي توفر نكهة إضافية. وبالنسبة إليّ فإن الأذن تحوّل الرغيف الجيد إلى رغيف رائع.

\begin{figure}[htb]
  \centering
  \includegraphics[width=\textwidth]{bread-scoring-angle}
  \caption[زاوية التشريط]{زاوية \LR{\ang{45}} التي تشريط بها العجين تكون بالنسبة إلى سطح العجين. وعند التشريط أكثر نحو الجانب عليك ضبط الزاوية لتحقيق الأذن في خبزك.}
  \label{fig:scoring-angle}
\end{figure}

ويُجرى الشق الفعلي بسكين حاد جدًا أو، والأفضل، بشفرة حلاقة. ويمكنك استخدام شفرة الحلاقة مباشرة أو تثبيتها على عصا خيزران. وتوفر شفرة الحلاقة مرونة أفضل من السكين الحادة. وفي كل الأحوال ينبغي أن تكون الشفرة حادة قدر الإمكان؛ فبذلك لن يتمزق العجين عند القطع بل سيحصل على شق نظيف غير ممزق.

ولتبسيط التشريط ينبغي أن يكون سطح عجينك مجففًا قليلًا؛ فبذلك يصبح إحداث الشق أسهل بكثير. ولهذا من الجوهري أن تفرك عجينك بقليل من الدقيق قبل وضعه في سلة التخمير؛ فالدقيق الجاف سيمتص بعض رطوبة الطبقات الخارجية من عجينك. وهذا مهم بدرجة خاصة عند العمل مع عجائن مخمَّرة نهائيًا في حرارة الغرفة. أما العجين المخمَّر نهائيًا باردًا فأسهل بكثير في التشريط بسبب لزوجته المنخفضة، بينما عجين حرارة الغرفة أصعب بكثير في التشريط ويمتزق شقه بسهولة أكبر. فمع شق ممزق لن ينتفخ العجين في الفرن كما ينبغي على الأرجح، ولن تحقق على الأرجح الأذن المذكورة سابقًا. ولهذا فتجفيف السطح مهم بدرجة خاصة، وسيصبح التشريط أسهل بكثير.

\begin{figure}[htb]
  \centering
  \includegraphics[width=\textwidth]{dry-dough-surface}
  \caption[تجفيف سطح العجين]{بتفريك الدقيق على سطح عجينك بعد التشكيل يجف الجزء الخارجي من العجين قليلًا، وهذا يجعل التشريط أسهل بكثير لأن الشق أقل عرضة للتمزق.}
  \label{fig:dried-out-dough-scoring}
\end{figure}

التشريط يتطلب تدربًا كثيرًا. ولهذا أنصح بالتدرب على إحداث الشق بعد بناء قوة العجين؛ فسيكون العجين مبلولًا ولزجًا جدًا. ويمكنك استخدام سكين حاد أو شفرة حلاقة للتدرب على التقنية. وانتظر بضع دقائق ثم دوّر العجين واستدره من جديد. ويمكنك التدرب ما شئت حتى ترضى عن تقنيتك. أما بعد التخمير النهائي فلن تكون لديك سوى فرصة واحدة للتدرب على التشريط؛ فهو إما إصابة وإما خطأ.

وحيلة إضافية يمكن أن تساعدك على الجمع بين فوائد التخمير النهائي في حرارة الغرفة وسهولة تشريط التخمير النهائي البارد هي وضع عجينك في المجمّد لمدة \LR{30} دقيقة قبل الخبز. وحالما تلاحظ أن عجينك كاد ينتهي من التخمير النهائي انقله إلى المجمّد؛ فسيجفف المجمّد سطح العجين مزيدًا مع خفض لزوجته في الوقت نفسه، مما يجعل التشريط أسهل.

وحيلة أخرى مثيرة للاهتمام أن تخبز عجينك \LR{30} ثانية بلا بخار؛ فالهواء الساخن سيجفف سطح العجين مزيدًا ويبسّط تقنية التشريط. وجرّب التوقيت لتحديد نقطة التوازن الشخصية الخاصة بك.
